\documentclass[11pt,a4paper]{article}
\usepackage[margin=24mm,headheight=14pt]{geometry}
\usepackage{fontspec}
\setmainfont{Latin Modern Roman}
\setsansfont{Latin Modern Sans}
\setmonofont{Latin Modern Mono}
\usepackage{babel}
\babelprovide[import,main]{portuguese}
\usepackage{amsmath,amsthm,mathtools}
\usepackage{unicode-math}
\setmathfont{Latin Modern Math}
\usepackage{microtype,graphicx,xcolor,booktabs,longtable,array,float}
\usepackage{enumitem}
\usepackage{tikz}
\usetikzlibrary{arrows.meta,positioning,calc,fit}
\usepackage[round,authoryear]{natbib}
\usepackage{caption}
\captionsetup{font=small,labelfont=bf,skip=7pt}
\usepackage{fancyhdr}
\usepackage{hyperref}
\definecolor{ink}{HTML}{172C3C}
\definecolor{accent}{HTML}{166F80}
\hypersetup{colorlinks=true,linkcolor=ink,citecolor=ink,urlcolor=accent,
 pdftitle={Biblioteca determinística de segmentos afins racionais em células diádicas: endereçamento, aproximação e codificação},
 pdfauthor={Victor Gabriel de Oliveira},
 pdfsubject={Construção racional, aproximação certificada e descrição aritmética}}
\pagestyle{fancy}
\fancyhf{}
\fancyhead[L]{\small\itshape Biblioteca determinística de segmentos afins racionais}
\fancyfoot[C]{\small\thepage}
\renewcommand{\headrulewidth}{0pt}
\setlength{\parindent}{1.25em}
\setlength{\parskip}{2pt}
\setlength{\emergencystretch}{2em}
\setlength{\textfloatsep}{15pt plus 3pt minus 3pt}
\setlength{\intextsep}{14pt plus 3pt minus 3pt}
\setlist{nosep,leftmargin=2em}
\newtheoremstyle{preciso}{8pt}{5pt}{\itshape}{}{\bfseries}{.}{.5em}{}
\theoremstyle{preciso}
\newtheorem{teorema}{Teorema}[section]
\newtheorem{lema}[teorema]{Lema}
\newtheorem{proposicao}[teorema]{Proposição}
\newtheorem{corolario}[teorema]{Corolário}
\theoremstyle{definition}
\newtheorem{definicao}[teorema]{Definição}
\newtheorem{algoritmo}[teorema]{Algoritmo}
\newtheorem{observacao}[teorema]{Observação}
\renewcommand{\proofname}{Demonstração}
\newcommand{\N}{\mathbb N}
\newcommand{\Z}{\mathbb Z}
\newcommand{\Q}{\mathbb Q}
\newcommand{\R}{\mathbb R}
\newcommand{\D}{\mathbb D}
\newcommand{\LL}{\mathcal L}
\newcommand{\Asm}{\mathcal A}
\newcommand{\FF}{\mathcal F}
\newcommand{\EE}{\mathcal E}
\newcommand{\RR}{\mathcal R}
\newcommand{\CC}{\mathcal C}
\newcommand{\heightq}{\operatorname{ht}}
\newcommand{\rank}{\operatorname{rank}}
\newcommand{\Eval}{\operatorname{Eval}}
\newcommand{\ser}{\operatorname{ser}}
\newcommand{\norm}[1]{\left\|#1\right\|_{\infty}}
\newcommand{\normK}[2]{\left\|#1\right\|_{\infty,#2}}
\newcommand{\eqn}[2]{\begin{equation}\label{eq:#1}#2\end{equation}}
\newcommand{\figura}[4]{\begin{figure}[!htbp]\centering
 \emph{Figura omitida nesta versao para formalizacao; a legenda original e preservada.}
 \caption{#3}\label{fig:#4}\end{figure}}
\title{\vspace{-18mm}\bfseries Biblioteca determinística de segmentos afins racionais em células diádicas\\[4pt]
 \large Endereçamento, aproximação e codificação}
\author{Victor Gabriel de Oliveira\\
 \small Pesquisador independente, Brasil\\[-1pt]
 \small\href{mailto:vg32518@gmail.com}{vg32518@gmail.com}}
\date{}
\begin{document}
\maketitle
\thispagestyle{plain}
\section*{Escopo da formalizacao integral dos resultados pendentes}
Este arquivo reune o texto matematico original das secoes 2 a 7, com os blocos auxiliares de entropia e catalogos incluidos diretamente. As secoes 2 a 6 documentam as dependencias ja formalizadas. Os alvos ainda pendentes sao os teoremas e proposicoes 4.3--4.5 e os resultados da secao 7, inclusive seus subsequentes blocos de erro, entropia e catalogos. Preservar definicoes, hipoteses e quantificadores de todos os enunciados; reportar separadamente qualquer resultado parcial ou condicional. O texto desta introducao e uma instrucao de escopo, nao um teorema do artigo.
\setcounter{section}{1}

% === INICIO DA FONTE: secoes/02_indexacao.tex ===
\section{Ferramentas de indexação}
Escreva $\N_0=\{0,1,2,\ldots\}$ e $T_m=m(m+1)/2$.

\subsection{Pareamento de Cantor}
Para $u,v\in\N_0$, defina
\eqn{01}{\pi(u,v)=T_{u+v}+v=\frac{(u+v)(u+v+1)}2+v.}
\begin{lema}[pareamento e cota]
$\pi:\N_0^2\to\N_0$ é bijetiva. Para $n\in\N_0$, sua inversa é
\eqn{02}{w=\left\lfloor\frac{\sqrt{8n+1}-1}{2}\right\rfloor,\qquad v=n-T_w,\qquad u=w-v.}
Além disso, $\pi(u,v)\le(u+v+1)^2$.
\end{lema}
\begin{proof}
Se $m=u+v$, então $T_m\le\pi(u,v)=T_m+v<T_{m+1}$, pois $0\le v\le m$. Assim, cada imagem pertence a um único bloco diagonal $[T_m,T_{m+1})$. Dado $n$, a expressão para $w$ em \eqref{eq:02} é precisamente o único inteiro com $T_w\le n<T_{w+1}$. Logo $0\le v=n-T_w\le w$, $u=w-v\ge0$ e $\pi(u,v)=n$, provando sobrejetividade. O bloco determina $w=u+v$ e, dentro dele, $v=n-T_w$; portanto o par é único. Por fim, com $x=u+v+1$, tem-se $T_{u+v}\le x^2/2$ e $v\le x^2/2$, donde $\pi(u,v)\le x^2$.
\end{proof}

\subsection{Inteiros e racionais por altura}
Defina $Z:\N_0\to\Z$ por
\eqn{03}{Z(t)=\begin{cases}t/2,&t\text{ par},\\-(t+1)/2,&t\text{ ímpar}.\end{cases}}
\begin{lema}
$Z$ é bijetiva, com $Z^{-1}(m)=2m$ para $m\ge0$ e $Z^{-1}(m)=-2m-1$ para $m<0$. Em particular, $Z^{-1}(m)\le2|m|+1$.
\end{lema}
\begin{proof}
Os índices pares percorrem $\N_0$ exatamente uma vez; os ímpares percorrem $\Z\setminus\N_0$ exatamente uma vez. As duas imagens são disjuntas e as fórmulas inversas resultam por resolução direta de \eqref{eq:03}.
\end{proof}

Todo $r\in\Q$ tem uma representação reduzida única $p/q$, com $p\in\Z$, $q\in\N$ e $\gcd(|p|,q)=1$. Adote a altura multiplicativa elementar
\eqn{04}{\heightq(p/q)=\max\{|p|,q\}.}
Ordene os pares reduzidos primeiro por altura, depois por $q$ e finalmente por $p$. Essa leitura define bijeções
\eqn{05}{Q:\N_0\to\Q,\qquad Q^*:\N_0\to\Q^*,}
na segunda das quais se omite $0/1$.

\begin{lema}[contagem por altura]
Seja $\Q_H=\{r\in\Q:\heightq(r)\le H\}$. Então
\eqn{06}{|\Q_H|\le3H^2\qquad\text{para todo }H\ge1,}
e existem $H_0$ e $c_1>0$ universais tais que
\eqn{07}{|\Q_H|\ge c_1H^2\qquad\text{para todo }H\ge H_0.}
Pode-se tomar $c_1=1/4$.
\end{lema}
\begin{proof}
A forma reduzida com denominador positivo é única: de $p/q=p'/q'$ e das duas coprimalidades segue $q\mid q'$ e $q'\mid q$, logo $q=q'$ e $p=p'$. Cada $r\in\Q_H$ corresponde então a um único par $(p,q)$ com $|p|\le H$ e $1\le q\le H$; há no máximo $(2H+1)H\le3H^2$ tais pares, o que prova \eqref{eq:06}.

Para \eqref{eq:07}, observe que $\Q_H$ contém todas as frações $p/q$ com $1\le q\le H$, $1\le p\le q$ e $\gcd(p,q)=1$, cuja quantidade é $\Phi(H)=\sum_{q\le H}\varphi(q)$. Pelo teorema clássico de Mertens sobre a soma de Euler \citep[Teorema 330]{hardy2008},
\[
\Phi(H)=\frac3{\pi^2}H^2+O(H\log H),
\]
e $3/\pi^2>0{,}30$. Logo $\Phi(H)\ge H^2/4$ para todo $H$ suficientemente grande.
\end{proof}

\begin{lema}[posto por altura]
As aplicações em \eqref{eq:05} são bijetivas. Se $\heightq(p/q)\le H$, então o posto de $p/q$ em qualquer uma delas é estritamente menor que $3H^2$.
\end{lema}
\begin{proof}
A ordenação esgota cada casca finita de altura antes de passar à seguinte; portanto lista todo racional uma única vez, o que dá a bijetividade. Todo elemento que precede $p/q$ pertence a $\Q_H$, e o lema~2.3 limita essa cardinalidade por $3H^2$. A remoção de $0/1$ em $Q^*$ só diminui a contagem.
\end{proof}

\subsection{Bijetividade e crescimento do índice}
A ordem de Calkin--Wilf fornece uma comparação entre bijetividade e crescimento do endereço. Seja $(S_n)_{n\ge0}$ a sequência diatômica de Stern,
\eqn{08}{S_0=0,\qquad S_1=1,\qquad S_{2n}=S_n,\qquad S_{2n+1}=S_n+S_{n+1},}
e $CW(n)=S_n/S_{n+1}$, $n\ge1$, que enumera $\Q_{>0}$ sem repetições.
\begin{lema}
Para todo $q\ge1$, $CW(2^{q-1})=1/q$.
\end{lema}
\begin{proof}
Da recorrência par, $S_{2^m}=S_1=1$. Também $S_{2^m+1}=m+1$: o caso $m=0$ é $S_2=1$ e
\[
S_{2^{m+1}+1}=S_{2(2^m)+1}=S_{2^m}+S_{2^m+1}=1+(m+1).
\]
Tomando $m=q-1$, resulta $CW(2^{q-1})=1/q$.
\end{proof}

O racional $1/q$ tem altura $q$, mas índice $2^{q-1}$; a escrita binária do índice exige $q$ bits, contra os $O(\log q)$ do lema~2.4. A ordenação por altura controla o endereço em função de $\heightq(r)$. A seção~7 determina o comprimento mínimo no pior caso por contagem e constrói uma enumeração que o realiza.
\figura{01_calkin_wilf}{.94}{Primeiros níveis da árvore de Calkin--Wilf. O ramo esquerdo contém $1/q$ no índice $2^{q-1}$; a ordenação por altura fornece índices polinomiais na altura.}{cw}

% === FIM DA FONTE: secoes/02_indexacao.tex ===


% === INICIO DA FONTE: secoes/03_biblioteca.tex ===
\section{Células diádicas e operador gerador}
\subsection{Células, posse e centros}
Para $k\in\N_0$ e $m\in\Z$, sejam
\eqn{09}{J_{k,m}=[m2^{-k},(m+1)2^{-k}],\qquad J_{k,m}^{\rm own}=[m2^{-k},(m+1)2^{-k}),}
e
\eqn{10}{x_0(k,m)=\left(m+\frac12\right)2^{-k}=\frac{2m+1}{2^{k+1}}.}
A célula fechada sustenta o segmento; a versão semiaberta atribui cada fronteira a um único proprietário. Escreva
\[
\D=\{j2^{-K}:j\in\Z,\ K\in\N_0\},\qquad \D_{\rm mid}=\D\setminus\Z.
\]
\begin{lema}[estrutura diádica]
Para cada $k$, as células $J_{k,m}^{\rm own}$ particionam $\R$ e
\eqn{11}{J_{k,m}=J_{k+1,2m}\cup J_{k+1,2m+1},\qquad |J_{k,m}|=2^{-k}.}
O mapa $(k,m)\mapsto x_0(k,m)$ é uma bijeção de $\N_0\times\Z$ sobre $\D_{\rm mid}$, que é denso em $\R$.
\end{lema}
\begin{proof}
Para $x\in\R$, o único índice de posse na escala $k$ é $m=\lfloor2^kx\rfloor$; isso prova a partição. A identidade de refinamento decorre de cindir $[2m2^{-(k+1)},(2m+2)2^{-(k+1)}]$ em seu ponto médio.

Pela \eqref{eq:10}, o centro tem numerador ímpar e denominador reduzido $2^{k+1}\ge2$, logo pertence a $\D_{\rm mid}$. Reciprocamente, cada $x\in\D_{\rm mid}$ possui forma reduzida única $j/2^K$, com $K\ge1$ e $j$ ímpar; tomando $k=K-1$ e $m=(j-1)/2$, obtém-se $x=x_0(k,m)$. A mesma unicidade prova a injetividade. Por fim, os centros de escala $k$ estão espaçados por $2^{-k}$; escolhendo o mais próximo de um real dado, a distância tende a zero com $k$, o que prova a densidade.
\end{proof}
\figura{02_celulas}{.92}{Células diádicas, centros e posse das fronteiras. Cada célula se divide em dois filhos; os centros determinam a célula de origem.}{celulas}

\subsection{Retas ancoradas e descompactação}
Dados $x_0,y_0\in\R$ e $s\ne0$, a reta ancorada é
\eqn{12}{\ell(x)=y_0+s(x-x_0)=a+bx,\qquad b=s,\quad a=y_0-sx_0.}
Para $n\in\N_0$, aplique sucessivamente a inversa do pareamento:
\eqn{13}{(u,v)=\pi^{-1}(n),\qquad(k,\mu)=\pi^{-1}(u),\qquad(r,t)=\pi^{-1}(v).}
\begin{lema}
O mapa $n\mapsto(k,\mu,r,t)$ de \eqref{eq:13} é uma bijeção de $\N_0$ sobre $\N_0^4$.
\end{lema}
\begin{proof}
Sua inversa explícita é
\[
(k,\mu,r,t)\longmapsto\pi\bigl(\pi(k,\mu),\pi(r,t)\bigr),
\]
uma composição de bijeções pelo lema~2.1.
\end{proof}
Defina
\eqn{14}{m=Z(\mu),\quad x_{0,n}=x_0(k,m),\quad y_{0,n}=Q(r),\quad s_n=Q^*(t),}
e converta para a forma global por
\eqn{15}{b_n=s_n,\qquad a_n=y_{0,n}-s_nx_{0,n}.}
Se $x_L=m2^{-k}$ e $x_R=(m+1)2^{-k}$, ponha
\eqn{16}{c_n=\min\{\ell_n(x_L),\ell_n(x_R)\},\qquad d_n=\max\{\ell_n(x_L),\ell_n(x_R)\}.}
\begin{definicao}[gerador e biblioteca]
Se $\mathcal I$ é o conjunto dos intervalos fechados não degenerados de $\R$, defina
\eqn{17}{\Gamma:\N_0\to\R\times(\R\setminus\{0\})\times\mathcal I,\qquad
\Gamma(n)=(a_n,b_n,[c_n,d_n]).}
A biblioteca é $\LL=\Gamma(\N_0)$. O átomo de endereço $n$ é a restrição de $\ell_n$ à célula $J_{k(n),m(n)}$.
\end{definicao}
Como $s_n\ne0$ e $x_L\ne x_R$, o intervalo em \eqref{eq:16} é não degenerado. A inclinação não nula torna a interface invertível no eixo espacial: uma peça constante teria imagem pontual e perderia a informação da célula. Para $\Gamma(n)=(a_n,b_n,[c_n,d_n])$, a pré-imagem espacial é
\eqn{18}{I_n^x=\left[
\min\left\{\frac{c_n-a_n}{b_n},\frac{d_n-a_n}{b_n}\right\},
\max\left\{\frac{c_n-a_n}{b_n},\frac{d_n-a_n}{b_n}\right\}\right].}

\subsection{Recuperação da célula e endereço canônico}
\begin{lema}[validade espacial exata]
Para todo $n\in\N_0$, o intervalo $I_n^x$ definido em \eqref{eq:18} satisfaz
\[
I_n^x=J_{k(n),m(n)}.
\]
\end{lema}
\begin{proof}
Como $b_n\ne0$, $\ell_n$ é uma bijeção afim estritamente monótona. Sua restrição a $[x_L,x_R]$ tem imagem $[c_n,d_n]$ por \eqref{eq:16}; logo $\ell_n^{-1}([c_n,d_n])=[x_L,x_R]$. Os dois quocientes em \eqref{eq:18} são precisamente $\ell_n^{-1}(c_n)$ e $\ell_n^{-1}(d_n)$, possivelmente em ordem invertida quando $b_n<0$.
\end{proof}
\begin{teorema}[unicidade do endereço]
$\Gamma$ é injetiva. Portanto, é uma bijeção de $\N_0$ sobre $\LL$ e cada átomo tem um único endereço.
\end{teorema}
\begin{proof}
Suponha $\Gamma(n)=\Gamma(n')=(a,b,[c,d])$. Pelo lema~3.4, a tripla recupera a célula comum $J=[x_L,x_R]$. Seu comprimento determina unicamente $k=-\log_2|J|$, e $m=2^kx_L$. Assim, também se recupera $x_0=(m+1/2)2^{-k}$. A inclinação local é $s=b$ e o valor ancorado é $y_0=a+bx_0$.

Pelas bijeções dos lemas~2.2, 2.4 e 3.1, os dados $(x_0,y_0,s)$ determinam unicamente $(k,\mu,r,t)$. O lema~3.2 determina então um único $n$. Logo $n=n'$.
\end{proof}
Em uma fronteira, duas células fechadas podem estar ativas. A célula semiaberta escolhe uma proprietária; em uma montagem contínua, as duas peças têm o mesmo valor no nó.
\begin{teorema}[dicionário de jatos racionais]
O mapa
\eqn{19}{n\longmapsto(x_{0,n},y_{0,n},s_n)}
é uma bijeção de $\N_0$ sobre $\D_{\rm mid}\times\Q\times\Q^*$. Sua imagem é densa em $\R^3$.
\end{teorema}
\begin{proof}
A bijetividade resulta da composição das bijeções em \eqref{eq:13}--\eqref{eq:14}. A densidade decorre da densidade de $\D_{\rm mid}$, de $\Q$ e de $\Q^*$ em $\R$, aplicada coordenada a coordenada na norma do máximo.
\end{proof}
O dicionário \eqref{eq:19} parametriza todos os 1-jatos com centro diádico não inteiro, valor racional e derivada racional não nula. Montagens contínuas ligam esses dados locais por valores de fronteira compartilhados.

% === FIM DA FONTE: secoes/03_biblioteca.tex ===


% === INICIO DA FONTE: secoes/04_montagens.tex ===
\section{Montagens e densidade em espaços de funções contínuas}
\begin{definicao}[montagem contínua]
Se $K=[\alpha,\beta]$ é um intervalo compacto não degenerado, $\Asm(K)$ é o conjunto das funções $g\in C(K)$ para as quais existe uma coleção finita $\Lambda$ de células diádicas, com interiores disjuntos e união contendo $K$, e endereços $n(J)$, $J\in\Lambda$, tais que
\eqn{20}{g|_{J\cap K}=\ell_{n(J)}|_{J\cap K},}
com valores coincidentes nas fronteiras compartilhadas.
\end{definicao}
Para $f\in C(K)$, denote o módulo de continuidade por
\eqn{21}{\omega(f;\delta)=\sup\{|f(x)-f(y)|:x,y\in K,\ |x-y|\le\delta\}.}
O resultado seguinte impõe simultaneamente nós diádicos, valores racionais e inclinações não nulas. A figura~\ref{fig:montagem} mostra a ligação das peças por nós comuns.
\figura{03_montagem_continua}{1}{Montagem contínua de átomos. À esquerda, cada cor identifica uma célula e seu segmento; peças consecutivas compartilham o mesmo valor no nó. À direita, o refinamento substitui células por filhos e aproxima a função-alvo com mais peças conectadas.}{montagem}

\begin{teorema}[densidade das montagens]
Para todo $f\in C[\alpha,\beta]$ e todo $\varepsilon>0$, existe $g\in\Asm([\alpha,\beta])$ tal que
\[
\normK{f-g}{[\alpha,\beta]}<\varepsilon.
\]
Todas as peças de $g$ pertencem a $\LL$ e têm inclinação não nula.
\end{teorema}
\begin{proof}
Estenda $f$ por constantes: $F(x)=f(\alpha)$ para $x<\alpha$, $F=f$ em $[\alpha,\beta]$ e $F(x)=f(\beta)$ para $x>\beta$. A extensão é uniformemente contínua. Escolha $k$ de modo que, com $h=2^{-k}$,
\eqn{22}{\omega(F;h)<\varepsilon/2.}
Tome os nós $t_j=jh$ desde $j_{\min}=\lfloor\alpha/h\rfloor$ até $j_{\max}=\lceil\beta/h\rceil$. Para cada $j$, escolha racional $v_j$ com $|v_j-F(t_j)|<\varepsilon/2$ e, sequencialmente, imponha $v_{j+1}\ne v_j$. Isso é possível porque cada intervalo aberto de escolha contém infinitos racionais e apenas um valor é proibido em cada passo.

Seja $g$ o interpolante linear dos pares $(t_j,v_j)$. Em $J_{k,j}$, sua inclinação é $(v_{j+1}-v_j)/h\in\Q^*$ e seu valor no centro é $(v_j+v_{j+1})/2\in\Q$. Pelo teorema~3.6, a peça é o único átomo com esses parâmetros. Os valores nodais compartilhados tornam $g$ contínua.

Se $P_kF$ é o interpolante dos valores exatos, então, para $x=(1-\theta)t_j+\theta t_{j+1}$,
\[
|F(x)-P_kF(x)|\le(1-\theta)|F(x)-F(t_j)|+\theta|F(x)-F(t_{j+1})|
\le\omega(F;h)<\varepsilon/2.
\]
Além disso, $P_kF-g$ é linear em cada célula e seus valores nodais têm módulo menor que $\varepsilon/2$, logo $\norm{P_kF-g}<\varepsilon/2$. A desigualdade triangular conclui.
\end{proof}
O número de células dessa cobertura uniforme é
\eqn{23}{N_k=\lceil\beta2^k\rceil-\lfloor\alpha2^k\rfloor=(\beta-\alpha)2^k+O(1).}

\subsection{Equivalência exata com redes ReLU racionais}
Escreva $\rho(z)=z_+=\max\{0,z\}$. A identidade em hinges abaixo conserva os nós diádicos, os parâmetros racionais e as inclinações não nulas da montagem.
\begin{teorema}[equivalência montagem--ReLU racional]
Seja $g\in\Asm(K)$ dada por $N$ folhas consecutivas, com nós $t_0<t_1<\cdots<t_N$ e inclinação racional não nula $s_j$ em $[t_j,t_{j+1}]\cap K$, $0\le j<N$. Então, em $K$,
\eqn{24}{g(x)=g(t_0)+s_0(x-t_0)+\sum_{j=1}^{N-1}(s_j-s_{j-1})\rho(x-t_j).}
Depois de expandido o primeiro termo afim, todos os coeficientes são racionais, todos os pontos de quebra internos são diádicos e \eqref{eq:24} é uma rede ReLU rasa normalizada com $N-1$ unidades hinge e um termo afim de saída. Unidades cujo peso de saída é zero podem ser removidas.

Reciprocamente, seja
\eqn{25}{R(x)=c+qx+\sum_{j=1}^Ma_j\rho(x-\tau_j),\qquad c,q,a_j\in\Q,\quad\tau_j\in\D,}
com nós distintos. Suponha que $K$ admita uma cobertura por células diádicas compatível com os $\tau_j$, que os valores nodais sejam racionais e que cada inclinação local
\eqn{26}{q+\sum_{\tau_j<x}a_j}
seja não nula. Então $R|_K\in\Asm(K)$. Uma malha diádica uniforme que refine todos os $\tau_j$ fornece essa montagem explicitamente.
\end{teorema}
\begin{proof}
Na primeira folha, o lado direito de \eqref{eq:24} tem valor $g(t_0)$ em $t_0$ e inclinação $s_0$. Ao atravessar $t_j$, a derivada ganha exatamente $s_j-s_{j-1}$; por indução, sua inclinação na folha $j$ é $s_j$. Como ambos os lados são contínuos e coincidem no primeiro nó, coincidem em todo $K$. Cada $s_j$ é racional. Os nós internos são fronteiras de células diádicas, e o intercepto do primeiro átomo é racional por \eqref{eq:15}; logo a forma expandida tem somente parâmetros racionais.

Para a recíproca, refine uma grade diádica até que cada $\tau_j$ seja uma fronteira. Em qualquer célula aberta da grade, \eqref{eq:25} é afim, com inclinação \eqref{eq:26}. O valor no centro da célula é racional, pois o centro, os nós e todos os coeficientes são racionais. A hipótese de inclinação não nula permite aplicar o teorema~3.6, que fornece um único átomo da biblioteca nessa célula. A continuidade de $R$ faz os valores coincidirem nas fronteiras. A união das peças é, portanto, uma montagem no sentido da definição~4.1.
\end{proof}
A forma normalizada \eqref{eq:25} já implica valores nodais racionais. Se a arquitetura não oferecer conexão afim direta, $x=\rho(x)-\rho(-x)$ realiza o termo linear com duas unidades adicionais e viés constante na saída.

\begin{proposicao}[parâmetros e altura]
Na forma normalizada, a conversão de uma montagem com $N$ folhas armazena $2N$ racionais: o intercepto e a inclinação iniciais, mais um nó e um salto de inclinação para cada uma das $N-1$ unidades. Se todos os átomos têm altura no máximo $H\ge2$, a altura máxima dos parâmetros produzidos satisfaz
\eqn{27}{H_{\rm ReLU}\le2^{H+1}H^2.}
\end{proposicao}
\begin{proof}
Há $N-1$ pares $(t_j,s_j-s_{j-1})$, além dos dois parâmetros afins. Pela definição de altura atômica, $\heightq(s_j)\le H$, os nós satisfazem $\heightq(t_j)\le2^H$, e a forma geral do lema~7.11 dá $\heightq(s_j-s_{j-1})\le2H^2$. Para o intercepto inicial, escreva $a=y_0-sx_0$. Como $\heightq(x_0)\le2^H$,
\[
\heightq(a)\le2\heightq(y_0)\heightq(sx_0)\le2^{H+1}H^2.
\]
Isso prova \eqref{eq:27}. Como os demais parâmetros são produtos ou somas dos mesmos dados, essa cota também garante que a conversão é efetiva em aritmética racional exata.
\end{proof}
\figura{04_relu}{1}{Uma montagem racional com quatro folhas e a rede ReLU rasa dada pela identidade \eqref{eq:24}. O gráfico e os pesos são calculados a partir dos mesmos dados racionais.}{relu}
A identidade funcional do teorema~4.3 conduz aos formatos da seção~7.3, que distinguem a transmissão estrutural dos diádicos de sua expansão em frações reduzidas.

\subsection{Relação com o 1-jato de Taylor}
Se $J_{k,m}$ tem centro $x_0$ e largura $h$, o 1-jato de $f\in C^2(J_{k,m})$ é
\[
T_{x_0}f(x)=f(x_0)+f'(x_0)(x-x_0).
\]
\begin{proposicao}[pertencimento e aderência]
Se $f(x_0)\in\Q$ e $f'(x_0)\in\Q^*$, então $T_{x_0}f$ é o átomo de $\LL$ associado à célula. Em geral, para todo $\delta>0$ existem $q_0\in\Q$ e $q_1\in\Q^*$ tais que $\ell_q(x)=q_0+q_1(x-x_0)$ satisfaz
\eqn{28}{\normK{T_{x_0}f-\ell_q}{J_{k,m}}<\delta,}
\eqn{29}{\normK{f-\ell_q}{J_{k,m}}\le\frac{h^2}8\normK{f''}{J_{k,m}}+\delta.}
\end{proposicao}
\begin{proof}
No primeiro caso, $(x_0,f(x_0),f'(x_0))$ pertence ao contradomínio do teorema~3.6. No caso geral, escolha $q_0$ com $|f(x_0)-q_0|<\delta/2$ e $q_1\ne0$ com $|f'(x_0)-q_1|<\delta/h$. Como $|x-x_0|\le h/2$, obtém-se \eqref{eq:28}. O resto integral de Taylor fornece
\[
|f(x)-T_{x_0}f(x)|=\left|\int_{x_0}^x(x-t)f''(t)\,dt\right|
\le\frac{|x-x_0|^2}2\norm{f''}\le\frac{h^2}8\norm{f''}.
\]
A desigualdade triangular prova \eqref{eq:29}.
\end{proof}
Jatos reais arbitrários pertencem, portanto, à aderência local da biblioteca racional.

% === FIM DA FONTE: secoes/04_montagens.tex ===


% === INICIO DA FONTE: secoes/05_taxas.tex ===
\section{Taxas uniformes e constante assintótica}
Para $f\in C^2[\alpha,\beta]$, use a extensão
\eqn{30}{F(x)=\begin{cases}
f(\alpha)+f'(\alpha)(x-\alpha)+\frac12f''(\alpha)(x-\alpha)^2,&x<\alpha,\\
f(x),&\alpha\le x\le\beta,\\
f(\beta)+f'(\beta)(x-\beta)+\frac12f''(\beta)(x-\beta)^2,&x>\beta.
\end{cases}}
Então $F\in C^2(\R)$ e $\normK{F''}{\R}=\normK{f''}{[\alpha,\beta]}$. Seja $P_kF$ o interpolante nos nós $j2^{-k}$.

\begin{teorema}[cota aguda e limite]
Ponha $h=2^{-k}$ e $M=\normK{f''}{[\alpha,\beta]}$. Então
\eqn{31}{\normK{f-P_kF}{[\alpha,\beta]}\le\frac{h^2M}8,}
e
\eqn{32}{\lim_{k\to\infty}\frac8{h^2}\normK{f-P_kF}{[\alpha,\beta]}=M.}
Além disso, para todo $\delta>0$ existe $g_{k,\delta}\in\Asm([\alpha,\beta])$ com
\eqn{33}{\norm{f-g_{k,\delta}}\le\frac{h^2M}8+\delta.}
\end{teorema}
\begin{proof}
Em uma célula $[x_L,x_R]$ de largura $h$, ponha $e=F-P_kF$. Para $x\in(x_L,x_R)$, a função
\[
\psi(t)=e(t)-\frac{(t-x_L)(t-x_R)}{(x-x_L)(x-x_R)}e(x)
\]
anula-se em $x_L,x,x_R$. Três aplicações de Rolle produzem $\xi_x\in(x_L,x_R)$ com $\psi''(\xi_x)=0$; portanto
\eqn{34}{F(x)-P_kF(x)=\frac12F''(\xi_x)(x-x_L)(x-x_R).}
O módulo do produto é no máximo $h^2/4$, com igualdade no centro, o que prova \eqref{eq:31}.

Se $M=0$, \eqref{eq:32} é imediata. Se $M>0$, escolha $x_*\in[\alpha,\beta]$ com $|f''(x_*)|=M$. Para cada $k$ suficientemente grande, escolha um centro de escala $k$, $z_k\in[\alpha,\beta]$, com $z_k\to x_*$. Avaliando \eqref{eq:34} em $z_k$,
\[
|F(z_k)-P_kF(z_k)|=\frac{h^2}8|F''(\xi_k)|,
\]
onde $\xi_k$ pertence à célula de $z_k$. O diâmetro da célula tende a zero, logo $\xi_k\to x_*$ e $|F''(\xi_k)|\to M$. Isso dá o limite inferior; \eqref{eq:31} fornece o superior.

Para a última afirmação, racionalize todos os valores nodais da cobertura finita com erro menor que $\delta$ e mantenha valores adjacentes distintos, como no teorema~4.2. O interpolante resultante é uma montagem e dista menos que $\delta$ de $P_kF$ em norma uniforme. A desigualdade triangular conclui.
\end{proof}

\begin{corolario}[preservação da constante]
Existe uma sequência $g_k\in\Asm([\alpha,\beta])$ tal que
\[
\lim_{k\to\infty}\frac8{h^2}\norm{f-g_k}=M.
\]
\end{corolario}
\begin{proof}
Na construção do teorema~5.1, escolha erros nodais $\delta_k=o(h^2)$. Então
\[
\left|\norm{f-g_k}-\norm{f-P_kF}\right|\le\norm{g_k-P_kF}<\delta_k.
\]
Divida por $h^2/8$ e aplique \eqref{eq:32}.
\end{proof}
Com $N_k$ de \eqref{eq:23}, $h=(\beta-\alpha)N_k^{-1}(1+O(N_k^{-1}))$ e, portanto,
\eqn{35}{\norm{f-P_kF}\le\frac{(\beta-\alpha)^2M}8N_k^{-2}\bigl(1+O(N_k^{-1})\bigr).}
Quando os extremos já são nós da malha, o termo de borda desaparece.
\begin{table}[!htbp]\centering
\caption{Erro exato do interpolante de $e^x$ em $[0,1]$.}
\begin{tabular}{rrrr}\toprule
$k$&erro exato&$2^{-2k}e/8$&razão\\\midrule
2&$1{,}8774\times10^{-2}$&$2{,}1237\times10^{-2}$&$0{,}8840$\\
4&$1{,}2866\times10^{-3}$&$1{,}3273\times10^{-3}$&$0{,}9693$\\
6&$8{,}2310\times10^{-5}$&$8{,}2955\times10^{-5}$&$0{,}9922$\\
8&$5{,}1746\times10^{-6}$&$5{,}1847\times10^{-6}$&$0{,}9980$\\
10&$3{,}2389\times10^{-7}$&$3{,}2404\times10^{-7}$&$0{,}9995$\\\bottomrule
\end{tabular}
\end{table}
\figura{05_taxa_exponencial}{.98}{Erro do interpolante de $e^x$ e razão entre erro e cota. A razão converge a $1$, conforme \eqref{eq:32}.}{taxa}

% === FIM DA FONTE: secoes/05_taxas.tex ===


% === INICIO DA FONTE: secoes/06_refinamento.tex ===
\section{Refinamento diádico adaptativo certificado}
O refinamento usa majorantes locais validados de $|F''|$ para controlar o erro em cada célula.

\subsection{Fechamento operacional por aritmética intervalar}
Seja $\EE$ a menor classe de expressões que contém $x$ e as constantes racionais e é fechada por soma, subtração, produto, recíproco, $\exp$, $\sin$, $\cos$ e composição. Um recíproco é admissível em intervalos nos quais seu denominador não se anula. Essa classe contém os polinômios, as funções racionais sem polos no domínio e as funções elementares usadas nos exemplos.

Composição significa substituição sintática: em $u\circ v$, cada folha $x$ da árvore de $u$ é substituída pela árvore de $v$. Toda expressão composta é, assim, uma árvore das operações-base.

Escreva $\mathbb I_{\Q}$ para os intervalos fechados com extremos racionais. A rotina
\eqn{36}{\operatorname{BoundSecondDerivative}(E,J;P,d)
=[\underline m_J,\overline m_J]\in\mathbb I_{\Q}}
usa precisão transcendental $P$ e profundidade de subdivisão $d$. Ela propaga simultaneamente intervalos para $(E,E',E'')$, isto é, 2-jatos intervalares. Por exemplo,
\begin{align*}
(uv)''&=u''v+2u'v'+uv'',\\
(u^{-1})''&=2(u')^2u^{-3}-u''u^{-2},\\
(e^u)''&=e^u\bigl((u')^2+u''\bigr),\\
(\sin u)''&=-\sin(u)(u')^2+\cos(u)u'',\\
(\cos u)''&=-\cos(u)(u')^2-\sin(u)u''.
\end{align*}
As operações racionais são exatas. Para $\exp$, a redução do argumento a $|z|\le1/2$ é seguida por uma soma de Taylor com resto majorado por uma série geométrica racional; para $\sin$ e $\cos$, usa-se o resto de Lagrange racional e a interseção com $[-1,1]$. Elevar $P$ e $d$ reduz a sobrestimação.

\begin{lema}[restos transcendentais com extremos racionais]
Se $I=[a,b]\in\mathbb I_{\Q}$, as rotinas descritas acima produzem intervalos de extremos racionais que contêm $\exp(I)$, $\sin(I)$ e $\cos(I)$. Mais explicitamente, para $0\le z\le1/2$ e $n\ge0$,
\eqn{37}{0\le e^z-\sum_{j=0}^n\frac{z^j}{j!}
\le\frac{z^{n+1}}{(n+1)!}\frac1{1-z/(n+2)}.}
Se $R=\max\{|a|,|b|\}$ e
\[
P_n^{\sin}(z)=\sum_{j=0}^n(-1)^j\frac{z^{2j+1}}{(2j+1)!},\qquad
P_n^{\cos}(z)=\sum_{j=0}^n(-1)^j\frac{z^{2j}}{(2j)!},
\]
então, para todo $z\in I$,
\eqn{38}{|\sin z-P_n^{\sin}(z)|\le\frac{R^{2n+2}}{(2n+2)!},\qquad
|\cos z-P_n^{\cos}(z)|\le\frac{R^{2n+1}}{(2n+1)!}.}
Avaliar os polinômios em $I$, inflar pelos respectivos restos e intersectar com $[-1,1]$ preserva a inclusão.
\end{lema}
\begin{proof}
Na série de $e^z$, o primeiro termo omitido é $t_{n+1}=z^{n+1}/(n+1)!$. Para $j\ge n+1$, a razão entre termos consecutivos é $z/(j+1)\le z/(n+2)<1$; a cauda é, portanto, majorada pela série geométrica de primeiro termo $t_{n+1}$ e razão $z/(n+2)$, o que dá \eqref{eq:37}. Para um argumento racional $x$ arbitrário, escolha $r$ com $z=|x|/2^r\le1/2$, aplique a inclusão anterior e eleve o intervalo ao quadrado $r$ vezes; se $x<0$, tome o recíproco do intervalo positivo obtido para $-x$. Em $I=[a,b]$, a monotonicidade de $e^x$ permite usar o extremo inferior da inclusão em $a$ e o superior da inclusão em $b$.

As duas desigualdades em \eqref{eq:38} são o resto de Lagrange dos polinômios de Maclaurin de graus $2n+1$ e $2n$, pois todas as derivadas de seno e cosseno têm módulo no máximo $1$. A avaliação intervalar contém a imagem de cada polinômio; inflar pelo raio do resto contém a função. Como seno e cosseno já pertencem a $[-1,1]$, a interseção não remove nenhum valor possível. Todas as operações, fatoriais e extremos envolvidos são racionais.
\end{proof}

\begin{proposicao}[correção do majorante intervalar]
Se $E\in\EE$ define uma função $C^2$ em $J$ e a rotina \eqref{eq:36} termina com sucesso, então
\eqn{39}{E''(x)\in[\underline m_J,\overline m_J]\ (x\in J),\qquad
M_{\rm int}(J):=\max\{|\underline m_J|,|\overline m_J|\}\ge\normK{E''}{J}.}
Se algum intervalo de denominador contém zero mesmo após a subdivisão permitida, a rotina devolve falha e nenhum número é aceito como certificado.
\end{proposicao}
\begin{proof}
Para constantes e para $x$, a inclusão é imediata. Soma, produto e recíproco preservam inclusão pelas operações elementares de intervalos; o recíproco só é executado quando zero está fora do intervalo do denominador. As fórmulas de 2-jatos acima e o lema~6.1 dão o passo de indução estrutural para toda expressão de $\EE$. Pela normalização sintática declarada antes do lema, uma composição genérica já foi expandida nessa árvore e não acrescenta um caso independente. Subdividir $J$ preserva a cobertura, e tomar o fecho convexo dos intervalos obtidos em suas partes contém a união das imagens de $E''$. Isso prova a primeira inclusão e, tomando módulos, \eqref{eq:39}. A regra de falha impede que uma divisão não validada produza uma falsa garantia.
\end{proof}
A propagação registra a largura dos intervalos, o número de caixas e a profundidade usados. Subdivisão e Taylor racional permitem apertar as inclusões.

\begin{algoritmo}[refinamento certificado]
Dados $f\in C^2[\alpha,\beta]$, tolerâncias $\varepsilon_I,\varepsilon_R>0$ e a extensão $F$ de \eqref{eq:30}, use, para toda célula $J$, um majorante validado
\eqn{40}{\normK{F''}{J}\le M(J)\le M_*<\infty.}
Quando $F$ é fornecida por uma expressão de $\EE$, tome $M(J)=M_{\rm int}(J)$ de \eqref{eq:39}; se a certificação falhar, aborte sem emitir uma montagem garantida. Para uma função $C^2$ fora dessa interface, \eqref{eq:40} é um contrato de entrada.
\begin{enumerate}[label=(\roman*),topsep=6pt,itemsep=3pt]
\item Inicie com uma cobertura diádica finita de $[\alpha,\beta]$.
\item Para cada folha $J$, calcule
\eqn{41}{\eta_{\rm cert}(J)=\frac{|J|^2M(J)}8.}
\item Enquanto existir folha com $\eta_{\rm cert}(J)>\varepsilon_I$, substitua-a por seus dois filhos.
\item Nos nós finais, escolha racionais compartilhados $v_t$ com $|v_t-F(t)|<\varepsilon_R$ e valores adjacentes distintos.
\item Emita, em cada folha, o endereço do átomo que interpola seus dois valores nodais.
\end{enumerate}
\end{algoritmo}

\begin{teorema}[correção, término e tamanho]
O algoritmo~6.3 termina e produz $g\in\Asm([\alpha,\beta])$ com
\eqn{42}{\norm{f-g}\le\varepsilon_I+\varepsilon_R.}
Se a cobertura inicial tem $R_0$ raízes de largura máxima $h_0$ e $M_*>0$, nenhuma ramificação ultrapassa
\eqn{43}{r_*=\max\left\{0,\left\lceil\frac12\log_2\left(\frac{h_0^2M_*}{8\varepsilon_I}\right)\right\rceil\right\}}
níveis, e há no máximo $R_02^{r_*}$ folhas. Se $M_*=0$, pode-se tomar $r_*=0$.
\end{teorema}
\begin{proof}
Após $r$ subdivisões, uma folha descendente de uma raiz tem $|J|\le h_02^{-r}$ e, por \eqref{eq:40},
\[
\eta_{\rm cert}(J)\le\frac{h_0^2M_*}8\,4^{-r}.
\]
O lado direito é no máximo $\varepsilon_I$ para $r\ge r_*$; logo nenhuma folha é subdividida além dessa profundidade. A árvore é finita e cada raiz gera no máximo $2^{r_*}$ folhas.

Na partição final, \eqref{eq:34} e \eqref{eq:40} implicam
\[
\normK{F-P_\Lambda F}{J}\le\frac{|J|^2}8\normK{F''}{J}\le\eta_{\rm cert}(J)\le\varepsilon_I.
\]
A racionalização nodal produz, pelos teoremas~3.5 e 3.6, uma montagem de átomos de endereço único e com $\norm{P_\Lambda F-g}<\varepsilon_R$. Como $F=f$ no intervalo original, a desigualdade triangular prova \eqref{eq:42}.
\end{proof}
\begin{figure}[H]\centering
\emph{Ilustracao omitida na versao de formalizacao.}
\caption{Refinamento certificado para $f(x)=[1+64(x-5/8)^2]^{-1}$. Células com certificado acima de $\varepsilon_I$ são subdivididas. A malha se concentra na região de maior curvatura, e o painel inferior compara os certificados finais com a tolerância. A faixa superior representa o limite uniforme $\varepsilon_I+\varepsilon_R$.}
\label{fig:refinamento}
\end{figure}

\begin{proposicao}[calibração exata para curvatura constante]
Suponha que $[\alpha,\beta]=J_0$ seja uma única célula diádica inicial, $F''\equiv\kappa\ne0$ em $J_0$ e $M(J)=|\kappa|$ em todo descendente. Pondo
\[
A=\frac{|J_0|^2|\kappa|}8,
\]
o número de folhas produzido com tolerância $\varepsilon_I$ é
\eqn{44}{N(\varepsilon_I)=2^{r(\varepsilon_I)},\qquad
r(\varepsilon_I)=\max\left\{0,\left\lceil\log_4\left(\frac A{\varepsilon_I}\right)\right\rceil\right\}.}
Consequentemente,
\eqn{45}{\begin{aligned}
\liminf_{\varepsilon_I\downarrow0}\sqrt{\varepsilon_I}\,N(\varepsilon_I)
&=\frac1{\sqrt8}\int_\alpha^\beta|f''(x)|^{1/2}\,dx,\\
\limsup_{\varepsilon_I\downarrow0}\sqrt{\varepsilon_I}\,N(\varepsilon_I)
&=\frac1{\sqrt2}\int_\alpha^\beta|f''(x)|^{1/2}\,dx.
\end{aligned}}
A quantização diádica introduz, nesse caso, um fator dois entre os limites de acumulação.
\end{proposicao}
\begin{proof}
Após $r$ subdivisões uniformes, cada folha tem largura $|J_0|2^{-r}$ e certificado $A4^{-r}$. Como todos os certificados de um mesmo nível são iguais, o algoritmo para no primeiro $r$ com $A4^{-r}\le\varepsilon_I$, o que prova \eqref{eq:44}. Para $A4^{-r}\le\varepsilon_I<A4^{-(r-1)}$, tem-se
\[
\sqrt A\le\sqrt{\varepsilon_I}\,2^r<2\sqrt A.
\]
Os extremos inferior e superior são atingidos, respectivamente, nos limiares $\varepsilon_I=A4^{-r}$ e por aproximação pela esquerda aos limiares $A4^{-(r-1)}$. Por fim, $\int_\alpha^\beta|f''|^{1/2}=|J_0|\sqrt{|\kappa|}$.
\end{proof}

\begin{observacao}[diferença segunda e erro no centro]
Se $x_0=(x_L+x_R)/2$, então
\[
|f(x_0)-P_Jf(x_0)|=\frac12|f(x_L)-2f(x_0)+f(x_R)|.
\]
Isso mede o erro no centro. Para $f(x)=x^3$ em $[-1,1]$, $P_Jf(x)=x$ e a diferença segunda central é zero, enquanto
\[
\norm{x^3-x}=\frac2{3\sqrt3}\approx0{,}3849.
\]
O controle em todo o intervalo é fornecido pelo majorante de \eqref{eq:40}.
\end{observacao}

% === FIM DA FONTE: secoes/06_refinamento.tex ===


% === INICIO DA FONTE: secoes/07_codificacao.tex ===
\section{Endereços, entropia, descrição e busca}
\subsection{Comprimento do índice canônico}
Para um átomo com parâmetros $(k,m,y_0,s)$, defina
\eqn{46}{H=\max\{2,k+1,|m|+1,\heightq(y_0),\heightq(s)\}.}
Escreva $\LL_H\subset\LL$ para o conjunto dos átomos de altura no máximo $H$.

\begin{teorema}[cota polinomial do endereço]
O único endereço $n$ do átomo satisfaz
\eqn{47}{n=O(H^8),\qquad\log_2(n+1)=O(\log H).}
As constantes são universais para a enumeração fixada.
\end{teorema}
\begin{proof}
Pelo lema~2.4, os postos $r$ de $y_0$ e $t$ de $s$ são $O(H^2)$. Pelo lema~2.2, $\mu=Z^{-1}(m)=O(H)$, e $k=O(H)$. A cota do lema~2.1 dá
\[
u=\pi(k,\mu)=O(H^2),\qquad v=\pi(r,t)=O(H^4),
\]
e, por uma última aplicação,
\[
n=\pi(u,v)=O\bigl((u+v+1)^2\bigr)=O(H^8).
\]
Como $H\ge2$, tomar logaritmos fornece a segunda afirmação.
\end{proof}

\subsection{Contagem por altura e consequência minimax para o endereço}
A contagem dos átomos determina o menor pior índice de uma enumeração; o aninhamento de Cantor determina o crescimento específico de $\Gamma$.
\begin{lema}[cardinalidade das cascas de altura]
Para todo $H\ge2$,
\eqn{48}{|\LL_H|\le18H^6.}
Além disso, existem uma constante universal $c_{\LL}>0$ e $H_1\ge2$ tais que, para todo $H\ge H_1$,
\eqn{49}{|\LL_H|\ge c_{\LL}H^6.}
Em particular, $|\LL_H|=\Theta(H^6)$.
\end{lema}
\begin{proof}
Pelo teorema~3.6, um átomo é determinado bijetivamente pela quádrupla $(k,m,y_0,s)\in\N_0\times\Z\times\Q\times\Q^*$, e a altura \eqref{eq:46} limita cada coordenada separadamente. Assim
\[
|\LL_H|=\#\{k\le H-1\}\cdot\#\{|m|\le H-1\}\cdot|\Q_H|\cdot|\Q_H\setminus\{0\}|,
\]
com $\Q_H$ como no lema~2.3. O primeiro fator é $H$, o segundo é $2H-1$, e os dois últimos estão entre $c_1H^2$ e $3H^2$ para $H$ grande. Logo
\begin{align*}
|\LL_H|&\le H\cdot2H\cdot3H^2\cdot3H^2=18H^6,\\
|\LL_H|&\ge H\cdot H\cdot c_1H^2\cdot(c_1H^2-1)\ge\frac{c_1^2}2H^6
\end{align*}
para todo $H$ suficientemente grande. Tome $c_{\LL}=c_1^2/2$.
\end{proof}

Com orçamentos desacoplados,
\[
\LL_{K_0,M_0,Y,S}:=\{(k,m,y_0,s):k\le K_0,\ |m|\le M_0,\ \heightq(y_0)\le Y,\ \heightq(s)\le S\},
\]
a bijeção do teorema~3.6 dá
\[
|\LL_{K_0,M_0,Y,S}|=(K_0+1)(2M_0+1)|\Q_Y|(|\Q_S|-1).
\]
A ordem é $(K_0+1)(M_0+1)Y^2S^2$ para $Y,S$ grandes. O calibre \eqref{eq:46} acopla $K_0=M_0=H-1$ e $Y=S=H$, produzindo $H^6$; já $K_0=M_0=\lfloor\log_2H\rfloor$ e $Y=S=H$ produziriam $H^4(\log H)^2$.

\begin{proposicao}[ordem exata do pior endereço aritmético]
Se
\[
M_\Gamma(H):=\max_{a\in\LL_H}\Gamma^{-1}(a),
\]
então
\eqn{50}{M_\Gamma(H)=\Theta(H^8).}
Consequentemente, o comprimento da maior palavra do gerador aritmético sobre $\LL_H$ é $8\log_2H+O(1)$ bits.
\end{proposicao}
\begin{proof}
A cota superior é o teorema~7.1. Para a inferior, escreva $N_H=|\Q_H|$. Pelo lema~2.3, $N_H\ge c_1H^2$ para $H$ suficientemente grande. Como as enumerações de \eqref{eq:05} listam por altura, existem $y_H\in\Q_H$ e $s_H\in\Q_H\setminus\{0\}$ cujos postos respectivos satisfazem
\[
r_H=N_H-1,\qquad t_H=N_H-2,
\]
e, portanto, $r_H,t_H\ge c_2H^2$ para alguma constante universal $c_2>0$. Considere o átomo de parâmetros $(k,m,y_0,s)=(0,0,y_H,s_H)$, que pertence a $\LL_H$. Como
\[
\pi(x,y)=\frac{(x+y)(x+y+1)}2+y\ge\frac{(x+y)^2}2,
\]
temos $v_H=\pi(r_H,t_H)\ge c_3H^4$. Além disso, $\pi(0,0)=0$, de modo que o endereço desse átomo satisfaz
\[
n_H=\pi\bigl(\pi(0,0),v_H\bigr)=\pi(0,v_H)\ge\frac{v_H^2}2\ge c_4H^8.
\]
Logo $M_\Gamma(H)=\Omega(H^8)$, que, junto da cota superior, prova \eqref{eq:50}. A afirmação em bits segue tomando logaritmos.
\end{proof}

\begin{teorema}[cota minimax relativa ao calibre]
Seja $A:\LL\to\N_0$ uma aplicação injetiva qualquer. Para todo $H\ge2$,
\eqn{51}{\max_{a\in\LL_H}A(a)\ge|\LL_H|-1.}
Consequentemente, para todo $H\ge H_1$, existe um átomo de altura no máximo $H$ com
\eqn{52}{A(\text{átomo})\ge c_{\LL}H^6-1,\qquad
\log_2\bigl(A(\text{átomo})+1\bigr)\ge6\log_2H-O(1).}
Nenhum endereçamento injetivo usa $o(\log H)$ bits no pior caso sobre $\LL_H$.
\end{teorema}
\begin{proof}
Por injetividade, $A(\LL_H)$ é um subconjunto de $\N_0$ com $|\LL_H|$ elementos distintos; seu máximo é, portanto, pelo menos $|\LL_H|-1$, o que prova \eqref{eq:51}. Aplique então \eqref{eq:49} e tome logaritmos.
\end{proof}
\begin{observacao}
O argumento vale para todo endereçamento injetivo, inclusive os definidos por tabelas, estado ou escolhas arbitrárias. Com a proposição~7.3, ele compara $8\log_2H+O(1)$ bits de $\Gamma$ a $6\log_2H+O(1)$ bits do limite minimax: a razão assintótica dos comprimentos é $4/3$.
\end{observacao}

\begin{definicao}[enumeração por altura da biblioteca]
Para $H\ge2$, a casca $\LL_H\setminus\LL_{H-1}$ é finita e não vazia. Ordene $\LL$ listando as cascas em ordem crescente de $H$ e, dentro de cada casca, ordenando as quádruplas $(k,m,y_0,s)$ lexicograficamente por
\[
\bigl(k,Z^{-1}(m),\rank_Q(y_0),\rank_{Q^*}(s)\bigr),
\]
onde $\rank$ denota o posto nas enumerações \eqref{eq:05}. A bijeção resultante é denotada $\Gamma^\sharp:\N_0\to\LL$, e $n^\sharp$ é o endereço de um átomo nessa ordem.
\end{definicao}

\begin{teorema}[endereço minimax exato para a exaustão]
A aplicação $\Gamma^\sharp$ é uma bijeção de $\N_0$ sobre $\LL$, computável nos dois sentidos, e para todo $H\ge2$
\eqn{53}{\max_{a\in\LL_H}n^\sharp(a)=|\LL_H|-1\le18H^6-1.}
Além disso, para qualquer endereçamento injetivo $A:\LL\to\N_0$,
\eqn{54}{\max_{a\in\LL_H}A(a)\ge|\LL_H|-1.}
Logo $\Gamma^\sharp$ realiza o mínimo possível do maior índice sobre $\LL_H$. Esse pior índice é $\Theta(H^6)$ e seu comprimento binário é minimax-ótimo para o calibre \eqref{eq:46}.
\end{teorema}
\begin{proof}
Cada casca é finita pelo lema~7.2 e a ordem dentro dela é total, de modo que a concatenação está bem definida e é bijetiva; a computabilidade nos dois sentidos segue de as cascas serem finitas e explicitamente enumeráveis. Por construção, $\Gamma^\sharp$ lista primeiro, e exatamente, os elementos de $\LL_H$; portanto $\Gamma^\sharp$ restrita a $\{0,\ldots,|\LL_H|-1\}$ é uma bijeção sobre $\LL_H$, o que dá a igualdade em \eqref{eq:53}, e \eqref{eq:48} dá a cota numérica. A desigualdade \eqref{eq:54} é \eqref{eq:51}: a imagem $A(\LL_H)$ tem $|\LL_H|$ elementos distintos de $\N_0$, logo seu máximo é ao menos $|\LL_H|-1$.
\end{proof}

\begin{observacao}[endereçamento e armazenamento]
$\Gamma$ descompacta o índice em célula, valor ancorado e inclinação pelo pareamento de Cantor. $\Gamma^\sharp$ ordena as mesmas quádruplas por altura e atinge o menor pior índice. A primeira usa três raízes quadradas inteiras na descompactação de \eqref{eq:13}; a segunda calcula cardinalidades de cascas, com somas de totientes truncadas e custo polinomial em $H$.

A conversão entre as duas ordens decodifica a quádrupla e a recodifica, preservando a célula e a função afim. A biblioteca tem, assim, a semântica fixada por $\Gamma$ e um codec minimax dado por $\Gamma^\sharp$.
\end{observacao}
Sob domínio, raiz e convenção de fronteira compartilhados, os endereços recuperam as células. O conjunto das folhas determina a árvore podada por fechamento ancestral até a raiz, e seus extremos determinam a ordem espacial. A transmissão usa apenas $N$, $H$ e a lista de endereços.

\begin{teorema}[serialização condicionada de montagens]
Fixe, como informação compartilhada, o domínio $K$, a raiz diádica inicial, a convenção de fronteira e a enumeração usada. Seja $g\in\Asm(K)$ descrita por uma árvore diádica com $N$ folhas, e seja $H\ge2$ a maior altura de seus átomos. Existe um fluxo de bits autodelimitado, decodificável sem conhecer previamente $N$ ou $H$, com comprimento
\eqn{55}{\begin{aligned}
L(g\mid K)&\le(2\lfloor\log_2N\rfloor+1)+(2\lfloor\log_2H\rfloor+1)
+N\left\lceil\log_2(18H^6)\right\rceil\\
&=N\bigl(6\log_2H+O(1)\bigr)+O(\log N+\log H)
\end{aligned}}
quando os átomos são transmitidos por $\Gamma^\sharp$. Com o gerador aritmético $\Gamma$, o termo por folha é $8\log_2H+O(1)$ bits.
\end{teorema}
\begin{proof}
Codifique $N$ e $H$, nessa ordem, pelo código gama de Elias, com comprimentos $2\lfloor\log_2N\rfloor+1$ e $2\lfloor\log_2H\rfloor+1$. Conhecido $H$, cada endereço de $\Gamma^\sharp$ ocupa exatamente $b_H=\lceil\log_2(18H^6)\rceil$ bits, pois $n^\sharp<18H^6$ pelo teorema~7.7. O cabeçalho $N$ informa quantos blocos de largura $b_H$ devem ser lidos. Os endereços recuperam os átomos e suas células; ordenar os extremos fornece a ordem espacial, e fechar as células por ancestrais recupera a árvore podada. O domínio compartilhado determina apenas o recorte nas extremidades. Isso prova a primeira cota e a decodificação sem esqueleto separado. Para $\Gamma$, use $\max_{a\in\LL_H}\Gamma^{-1}(a)=\Theta(H^8)$ da proposição~7.3.
\end{proof}
\begin{observacao}
Para extremos reais arbitrários, o domínio $K$ pode exigir informação infinita. Por isso, o teorema~7.9 mede a descrição condicionada ao domínio e à convenção de fronteira declarados.
\end{observacao}

\subsection{Três formatos prefix-free e propagação de altura}
Para $n\ge1$, seja
\[
\gamma(n)=0^{\lfloor\log_2n\rfloor}\operatorname{bin}(n)
\]
o código gama de Elias, e ponha
\[
\zeta(z)=\begin{cases}2z,&z\ge0,\\-2z-1,&z<0.\end{cases}
\]
Para $r=p/q$ reduzido, $q>0$, e para o diádico escrito canonicamente como $t=u2^{-k}$, com $u$ ímpar quando $k>0$, defina
\eqn{56}{C_{\Z}(z)=\gamma(\zeta(z)+1),}
\eqn{57}{C_{\rm rat}(p/q)=C_{\Z}(p)\gamma(q),\qquad L_{\rm rat}(p/q)=|C_{\rm rat}(p/q)|,}
\eqn{58}{C_{\rm dyad}(u,k)=C_{\Z}(u)\gamma(k+1),\qquad L_{\rm dyad}(u,k)=|C_{\rm dyad}(u,k)|.}
Cada componente é prefix-free e o cabeçalho determina a quantidade de componentes, tornando as concatenações decodificáveis sem separadores.

\begin{lema}[cálculo elementar de alturas]
Para quaisquer $r,s\in\Q$,
\eqn{59}{\heightq(r\pm s)\le2\heightq(r)\heightq(s),\qquad\heightq(rs)\le\heightq(r)\heightq(s).}
Em particular, se ambos têm altura no máximo $H$, as cotas respectivas são $2H^2$ e $H^2$. Para uma montagem de altura atômica máxima $H$, todo nó interno $t_j=u_j2^{-k_j}$ satisfaz $|u_j|\le H$ e $k_j\le H$ em sua forma reduzida; além disso,
\eqn{60}{\heightq(s_j-s_{j-1})\le2H^2,}
\eqn{61}{\heightq(s_jt_j)\le H2^H,}
\eqn{62}{\heightq\bigl(g(t_0)-s_0t_0\bigr)\le2^{H+1}H^2.}
\end{lema}
\begin{proof}
Escreva $r=p/q$ e $s=a/b$ em forma reduzida e ponha $R=\heightq(r)$, $S=\heightq(s)$. Antes de qualquer cancelamento, $r\pm s=(pb\pm aq)/(qb)$ e $rs=pa/(qb)$; os numeradores têm módulos no máximo $2RS$ e $RS$, respectivamente, e os denominadores são no máximo $RS$. Reduzir não aumenta altura, o que prova \eqref{eq:59}. Um nó é fronteira de uma célula $J_{k,m}$, logo sua forma não reduzida tem expoente $k\le H-1$ e numerador $m$ ou $m+1$, de módulo no máximo $H$; redução só diminui esses dados. As cotas \eqref{eq:60} e \eqref{eq:61} seguem da primeira parte e de $\heightq(t_j)\le2^H$. Finalmente, na primeira folha, $g(t_0)-s_0t_0=y_0-s_0x_0$; como $\heightq(x_0)\le2^H$, duas aplicações das cotas de produto e soma dão \eqref{eq:62}.
\end{proof}
Fixe a ordem espacial das folhas. Os três fluxos executáveis são:
\begin{enumerate}[label=(\roman*),topsep=5pt,itemsep=4pt]
\item $C_{\rm atom}(g)$: $\gamma(N)$, seguido, para cada folha, de $\gamma(k+1)C_{\Z}(m)C_{\rm rat}(y_0)C_{\rm rat}(s)$;
\item $C_{\rm ReLU,str}(g)$: $\gamma(N)$, o primeiro átomo na forma ancorada $C_{\rm dyad}(x_0)C_{\rm rat}(y_0)C_{\rm rat}(s_0)$, e os pares $C_{\rm dyad}(t_j)C_{\rm rat}(s_j-s_{j-1})$;
\item $C_{\rm ReLU,rat}(g)$: $\gamma(N)$ e a forma expandida $C_{\rm rat}(c)C_{\rm rat}(s_0)$, seguida dos pares $C_{\rm rat}(t_j)C_{\rm rat}(s_j-s_{j-1})$, onde $c=g(t_0)-s_0t_0$.
\end{enumerate}
Denote os comprimentos por $L_{\rm atom}$, $L_{\rm ReLU,str}$ e $L_{\rm ReLU,rat}$. O primeiro fluxo é uma alternativa estrutural à serialização por $\Gamma^\sharp$ do teorema~7.9; ambos têm a mesma ordem de comprimento.

\begin{teorema}[comparação rigorosa dos formatos]
Os três fluxos são prefix-free e possuem decodificadores exatos. Para toda montagem com $N$ folhas e altura atômica máxima $H\ge2$,
\eqn{63}{L_{\rm atom}(g)=O(N\log H),\quad L_{\rm ReLU,str}(g)=O(N\log H),\quad
L_{\rm ReLU,rat}(g)=O(NH),}
com constantes universais. A última ordem é atingida conjuntamente em $N$ e $H$: para todo $H\ge2$ e todo $N$ par com $2\le N\le H$, considere as folhas $J_{H-1,j}$, $0\le j<N$, com valores centrais zero e inclinações $(-1)^j$. Essa montagem contínua tem altura $H$, e
\eqn{64}{L_{\rm atom}=O(N\log H),\quad L_{\rm ReLU,str}=O(N\log H),\quad
L_{\rm ReLU,rat}=\Omega(NH).}
A equivalência funcional montagem--ReLU preserva, portanto, a ordem de bits no formato estrutural fixado.
\end{teorema}
\begin{proof}
O código gama é prefix-free; o mapa $\zeta$ é bijetivo entre $\Z$ e $\N_0$, e frações reduzidas e pares diádicos canônicos são únicos. Lendo o cabeçalho $N$, cada decodificador sabe quantos blocos seguintes consumir. No fluxo atômico, os blocos recuperam exatamente $(k,m,y_0,s)$ de cada folha. No fluxo ReLU estrutural, o primeiro bloco recupera $(x_0,y_0,s_0)$ e, portanto, $c=y_0-s_0x_0$; os demais recuperam os nós e saltos de inclinação da equação \eqref{eq:24}; o terceiro fluxo já transmite esses parâmetros expandidos. Logo os decodificadores são exatos.

Se $|z|\le A$, então $|C_{\Z}(z)|=O(\log(1+A))$, e se $\heightq(r)\le A$, então $L_{\rm rat}(r)=O(\log(1+A))$. As coordenadas dos átomos têm tamanho $O(H)$; os pares diádicos estruturais têm numerador e expoente $O(H)$; e os saltos têm altura $O(H^2)$ pelo lema~7.11. Isso prova as duas primeiras cotas de \eqref{eq:63}. Na forma ingênua, cada nó pode ter denominador $2^H$, e o intercepto pode ter altura $2^{H+1}H^2$; o lema~7.11 dá $O(H)$ bits por bloco e prova a terceira cota.

Na família final, as duas peças adjacentes ao nó $t_j=j2^{-(H-1)}$ assumem ali o mesmo valor $(-1)^{j-1}2^{-H}$, logo a montagem é contínua. Como $j<N\le H$, sua altura atômica máxima é exatamente $H$. Há ao menos $(N-2)/2$ índices internos ímpares $j$; para eles, $t_j$ já está reduzido e tem denominador $2^{H-1}$. Cada bloco racional correspondente custa $\Omega(H)$ bits por \eqref{eq:57}, e os saltos de inclinação são não nulos. Portanto $L_{\rm ReLU,rat}=\Omega(NH)$; os blocos atômicos e os pares diádicos estruturais custam $O(\log H)$ cada, dando as outras duas cotas. Para $N=2$, o intercepto $c=-2^{-H}$ fornece diretamente a mesma cota $\Omega(H)$. Isso prova \eqref{eq:64}.
\end{proof}

% === INICIO DA FONTE: secoes/07b_erro_entropia.tex ===
\subsection{Erro conjunto sob orçamentos de folhas e altura}
Os dois orçamentos abaixo são impostos à mesma representação por folhas.
\begin{definicao}[erro conjunto número--altura]
Para inteiros $N,H\ge2$, seja $\Asm_{N,H}(K)$ o conjunto das funções $g\in\Asm(K)$ que admitem uma representação como na definição~4.1 com no máximo $N$ folhas e com todos os átomos em $\LL_H$. Defina
\eqn{65}{E^{\LL}_{N,H}(f)_\infty:=\inf_{g\in\Asm_{N,H}(K)}\normK{f-g}{K},\qquad\inf\varnothing:=+\infty.}
Quando o conjunto abaixo não é vazio, ponha ainda
\eqn{66}{H^{\LL}_{\varepsilon,N}(f):=\min\{H\ge2:E^{\LL}_{N,H}(f)_\infty\le\varepsilon\}.}
\end{definicao}
$N$ mede a complexidade geométrica e $H$ controla resolução, posição e complexidade aritmética sob o calibre \eqref{eq:46}.

\begin{proposicao}[monotonicidade dos dois orçamentos]
Para todo $f\in C(K)$ e $N,H\ge2$,
\eqn{67}{E^{\LL}_{N+1,H}(f)_\infty\le E^{\LL}_{N,H}(f)_\infty,\qquad
E^{\LL}_{N,H+1}(f)_\infty\le E^{\LL}_{N,H}(f)_\infty.}
Se
\[
E^{\LL}_N(f)_\infty:=\inf\{\normK{f-g}{K}:g\in\Asm(K)\text{ admite no máximo }N\text{ folhas}\},
\]
então $E^{\LL}_{N,H}(f)_\infty\downarrow E^{\LL}_N(f)_\infty$ quando $H\to\infty$.
\end{proposicao}
\begin{proof}
As inclusões $\Asm_{N,H}(K)\subseteq\Asm_{N+1,H}(K)$ e $\Asm_{N,H}(K)\subseteq\Asm_{N,H+1}(K)$ provam \eqref{eq:67}. Toda representação finita tem altura máxima finita, logo a união de $\Asm_{N,H}(K)$ sobre $H\ge2$ é exatamente a classe usada para definir $E^{\LL}_N$. O ínfimo sobre essa união é o limite da sequência decrescente de ínfimos.
\end{proof}
Fixe $R,M>0$ e considere
\eqn{68}{\FF_{R,M}=\{f\in C^2[0,1]:\norm f\le R,\ \norm{f''}\le M\}.}

\begin{lema}[controle uniforme da primeira derivada]
Se $f\in\FF_{R,M}$ e
\eqn{69}{C_0:=4R+\frac M2,}
então $\norm{f'}\le C_0$.
\end{lema}
\begin{proof}
Para $x\in[0,1]$, escolha $a\in\{0,1\}$ com $|a-x|\ge1/2$. Taylor com resto de Lagrange dá
\eqn{70}{|f'(x)|\le\frac{|f(a)|+|f(x)|}{|a-x|}+\frac M2|a-x|
\le4R+\frac M2=C_0.}
\end{proof}

\begin{teorema}[controle simultâneo de erro, folhas, altura e bits]
Defina
\eqn{71}{C_H:=\max\{4(R+1),2(C_0+1)\}.}
Para toda potência diádica $n=2^r\ge2$ e todo $f\in\FF_{R,M}$, existe uma montagem $g_n$ com exatamente $n$ folhas tal que
\eqn{72}{\norm{f-g_n}\le\frac{M+4}{8n^2},}
\eqn{73}{\heightq_{\max}(g_n)\le C_Hn^2.}
Nos três formatos do teorema~7.12, a mesma montagem satisfaz
\eqn{74}{L_{\rm atom}(g_n),L_{\rm ReLU,str}(g_n),L_{\rm ReLU,rat}(g_n)
=O_{R,M}(n\log(2n)).}
Consequentemente, para todo inteiro $N\ge2$,
\eqn{75}{E^{\LL}_{N,\lceil C_HN^2\rceil}(f)_\infty<\frac{M+4}{2N^2}.}
A construção é determinística: quando $f(t_j)$ é equidistante dos dois pontos admissíveis da grade alternada, escolhe-se o ponto menor, ou, equivalentemente no denominador comum de \eqref{eq:77}, o menor numerador $b_j$.
\end{teorema}
\begin{proof}
Nos nós $t_j=j/n$, $0\le j\le n$, ponha $\Delta=n^{-2}$ e use as grades alternadas
\eqn{76}{G_j^{(n)}=\Delta\left(\Z+\frac{j\bmod2}{2}\right).}
Escolha $v_j\in G_j^{(n)}$ mais próximo de $f(t_j)$; no empate, escolha o ponto menor. Então $|v_j-f(t_j)|\le\Delta/2$, e grades consecutivas são disjuntas, de modo que $v_{j+1}\ne v_j$. O interpolante $g_n$ desses dados tem, portanto, $n$ peças com inclinações racionais não nulas; cada peça é um átomo pelo teorema~3.6. A estimativa de interpolação e a estabilidade convexa nos valores nodais dão
\[
\norm{f-g_n}\le\frac{M}{8n^2}+\frac1{2n^2}=\frac{M+4}{8n^2},
\]
o que prova \eqref{eq:72}.

Resta acompanhar a aritmética sem recorrer a uma cota genérica de soma de racionais. Escreva, de modo comum,
\eqn{77}{v_j=\frac{b_j}{2n^2},\qquad b_j\in\Z,\quad b_j\equiv j\pmod2.}
Como $|v_j|\le R+(2n^2)^{-1}$,
\eqn{78}{|b_j|\le2Rn^2+1.}
Na folha $J_{r,j}$, o valor central e a inclinação são
\eqn{79}{y_{0,j}=\frac{b_j+b_{j+1}}{4n^2},\qquad s_j=\frac{b_{j+1}-b_j}{2n}.}
O numerador de $s_j$ é ímpar, portanto a segunda fração já está reduzida e é não nula. Pelo lema~7.15,
\eqn{80}{|b_{j+1}-b_j|=2n^2|v_{j+1}-v_j|
\le2n^2\left(\frac{C_0}n+\frac1{n^2}\right)=2C_0n+2.}
Reduzir $y_{0,j}$ só diminui sua altura; logo
\eqn{81}{\heightq(y_{0,j})\le4(R+1)n^2,\qquad\heightq(s_j)\le2(C_0+1)n.}
Além disso, $r+1\le n$ e $|j|+1\le n$. A definição \eqref{eq:46} e \eqref{eq:71} provam \eqref{eq:73}.

Para os bits, no formato ReLU expandido o intercepto é $c=v_0$, $\heightq(c)\le2(R+1)n^2$; cada nó interno tem altura no máximo $n$; e, de \eqref{eq:80},
\[
\heightq(s_j-s_{j-1})\le4(C_0+1)n.
\]
Assim, tanto os parâmetros atômicos quanto os parâmetros ReLU desta construção têm altura polinomial em $n$. Nos códigos concretos \eqref{eq:57}--\eqref{eq:58}, cada bloco usa $O_{R,M}(\log(2n))$ bits; há $O(n)$ blocos em qualquer dos três fluxos. Isso prova \eqref{eq:74}. Observe que a cota geral $O(NH)$ do formato racional ingênuo não é saturada por esta família: seus nós têm denominador apenas $n$, e não $2^H$.

Por fim, dado $N\ge2$, tome $n=2^{\lfloor\log_2N\rfloor}$. Então $N/2<n\le N$. A montagem já construída usa no máximo $N$ folhas, tem altura no máximo $C_Hn^2\le C_HN^2$ e erro menor que $(M+4)/(2N^2)$, provando \eqref{eq:75}.
\end{proof}

\begin{corolario}[altura suficiente para precisão prescrita]
Para $\varepsilon>0$, seja
\eqn{82}{n_\varepsilon=
2^{\left\lceil\log_2\max\left\{2,\sqrt{\frac{M+4}{8\varepsilon}}\right\}\right\rceil},
\qquad H_\varepsilon=\lceil C_Hn_\varepsilon^2\rceil.}
Então, uniformemente para $f\in\FF_{R,M}$,
\eqn{83}{E^{\LL}_{n_\varepsilon,H_\varepsilon}(f)_\infty\le\varepsilon,}
\eqn{84}{n_\varepsilon=O_{R,M}(\varepsilon^{-1/2}),\qquad
H_\varepsilon=O_{R,M}(\varepsilon^{-1}),}
\eqn{85}{L_{\rm atom},L_{\rm ReLU,str},L_{\rm ReLU,rat}
=O_{R,M}\left(\varepsilon^{-1/2}\log\frac2\varepsilon\right).}
Em particular, $H^{\LL}_{\varepsilon,n_\varepsilon}(f)\le H_\varepsilon$.
\end{corolario}
\begin{proof}
A escolha de $n_\varepsilon$ torna o lado direito de \eqref{eq:72} no máximo $\varepsilon$. A menor potência de dois acima de um número positivo é menor que duas vezes esse número, salvo o piso fixo $2$; as três ordens seguem então do teorema~7.16.
\end{proof}
O denominador diádico comum de \eqref{eq:77} e o controle das diferenças adjacentes fornecem a cota construtiva $H_\varepsilon=O(\varepsilon^{-1})$.

A tabela~\ref{tab:bits} reúne erro, altura e descrição para $f_*(x)=4x(x-1)$, $R=1$ e $M=8$. Nesse caso, $C_0=8$, $C_H=18$, o erro exato é $49/(64N^2)$ e a altura efetiva é $4N^2$. Os três comprimentos são contados bit a bit pelos codificadores.
\begin{table}[!htbp]\centering
\caption{Construção do teorema~7.16 para $f_*(x)=4x(x-1)$. $e_{\rm cert}$ é a cota \eqref{eq:72}, $e_{\rm exato}$ é obtido nos extremos e no ponto crítico de cada folha, e $H_N=18N^2$.}\label{tab:bits}
\small\renewcommand{\arraystretch}{1.4}
\begin{tabular}{rrrrrrrr}\toprule
$N$&$e_{\rm cert}$&$e_{\rm exato}$&$\heightq_{\max}$&$H_N$&$L_{\rm atom}$&$L_{\rm str}$&$L_{\rm rat}$\\\midrule
2&$3/8$&$49/256$&16&72&77&59&37\\
4&$3/32$&$49/1024$&64&288&195&115&89\\
8&$3/128$&$49/4096$&256&1152&513&235&211\\
16&$3/512$&$49/16384$&1024&4608&1243&495&497\\
32&$3/2048$&$49/65536$&4096&18432&2921&1067&1159\\
64&$3/8192$&$49/262144$&16384&73728&6731&2329&2669\\\bottomrule
\end{tabular}
\end{table}

\subsection{Entropia de uma classe limitada de \texorpdfstring{$C^2$}{C2}}
Se $\mathcal N_\varepsilon(\FF_{R,M})$ é o menor número de bolas de raio $\varepsilon$ em norma uniforme que cobre a classe, escreva
\[
\mathcal H_\varepsilon(\FF_{R,M})=\log_2\mathcal N_\varepsilon(\FF_{R,M}).
\]
\begin{teorema}[entropia e codebook explícito]
Existem $\varepsilon_0,c,C,C_1>0$, dependentes apenas de $R$ e $M$, tais que, para $0<\varepsilon<\varepsilon_0$,
\eqn{86}{c\varepsilon^{-1/2}\le\mathcal H_\varepsilon(\FF_{R,M})\le C\varepsilon^{-1/2}.}
Existe também um codebook finito $\CC_\varepsilon\subset\Asm([0,1])$, composto somente por montagens de átomos de $\LL$, com distorção no máximo $\varepsilon$. Quando o codebook e seu decodificador são informação compartilhada, o comprimento de seu índice global de comprimento fixo é
\eqn{87}{L_\varepsilon^{\rm global}(\CC_\varepsilon):=
\lceil\log_2|\CC_\varepsilon|\rceil\le C_1\varepsilon^{-1/2}.}
Esse codebook atinge em ordem a entropia métrica da classe em cardinalidade e em bits de índice global.
\end{teorema}
\begin{proof}
\emph{Cota superior em \eqref{eq:86}.} Tome
\[
N=\left\lceil\sqrt{M/\varepsilon}\right\rceil+1,\qquad h=1/N,
\]
e seja $P_Nf$ o interpolante nos nós $j/N$. A estimativa local \eqref{eq:34}, válida para intervalos arbitrários, fornece $\norm{f-P_Nf}\le Mh^2/8\le\varepsilon/8$. Quantize cada $v_j=f(j/N)$ no múltiplo mais próximo de $\Delta=\varepsilon/4$, obtendo $v'_j$ com $|v'_j-v_j|\le\varepsilon/8$. O interpolante quantizado $g$ satisfaz $\norm{f-g}\le\varepsilon/4<\varepsilon$.

Resta contar os vetores quantizados que podem ocorrer. Há $O(\varepsilon^{-1})$ escolhas para $v'_0$. Dado $v'_0$, \eqref{eq:70} implica
\[
|v'_1-v'_0|\le C_0h+\varepsilon/4,
\]
e portanto há $O((N\varepsilon)^{-1})=O(\varepsilon^{-1/2})$ escolhas para $v'_1$. Para $1\le j\le N-1$, a identidade integral
\[
v_{j+1}-2v_j+v_{j-1}=\int_0^h\int_0^h
f''((j-1)h+s+t)\,ds\,dt
\]
dá módulo no máximo $Mh^2\le\varepsilon$. A quantização altera a diferença segunda em no máximo $4(\varepsilon/8)=\varepsilon/2$. Dados $v'_{j-1}$ e $v'_j$, o próximo valor está, pois, em um intervalo de comprimento $3\varepsilon$, que contém no máximo uma constante universal $K$ de pontos da grade de passo $\varepsilon/4$. Consequentemente, o número total de vetores é
\[
O(\varepsilon^{-1})\,O(\varepsilon^{-1/2})\,K^N
=\exp\bigl(O(\varepsilon^{-1/2})\bigr),
\]
e a cota superior segue ao tomar $\log_2$.

\emph{Cota inferior em \eqref{eq:86}.} Fixe $\psi\in C_c^2(0,1)$ com $\norm\psi=1$. Ponha $A=4\varepsilon$ e
\[
w=\sqrt{\frac{A\norm{\psi''}}M}.
\]
Para $\varepsilon$ pequeno, $A\le R$ e cabem $m=\lfloor1/w\rfloor\asymp\varepsilon^{-1/2}$ intervalos consecutivos de largura $w$ em $[0,1]$. Se $a_j=(j-1)w$, defina
\[
\varphi_j(x)=A\psi\left(\frac{x-a_j}w\right),\qquad
f_S=\sum_{j\in S}\varphi_j,\quad S\subseteq\{1,\ldots,m\}.
\]
Os suportes são disjuntos, $\norm{f_S}\le A\le R$ e $\norm{f_S''}\le M$. Se $S\ne T$, então $\norm{f_S-f_T}=A=4\varepsilon$. Logo os $2^m$ elementos formam um empacotamento $4\varepsilon$-separado; uma bola de raio $\varepsilon$ contém no máximo um deles. Assim, $\mathcal N_\varepsilon\ge2^m$ e $\mathcal H_\varepsilon\ge m\ge c\varepsilon^{-1/2}$.

\emph{Codebook de montagens.} Seja $N=2^k$ a menor potência de dois com $M/(8N^2)\le\varepsilon/2$. Pela minimalidade da potência de dois, $N=\Theta(\varepsilon^{-1/2})$ e $h=N^{-1}=\Theta(\varepsilon^{1/2})$. Defina
\[
\delta=2^{-\lceil\log_2(1/\varepsilon)\rceil},\qquad
\varepsilon/2<\delta\le\varepsilon,
\]
e, no nó $t_j=j/N$, use a grade finita alternada
\eqn{88}{G_j=\delta\left(\Z+\frac{j\bmod2}{2}\right)\cap[-R-\delta,R+\delta].}
Para cada $f$, escolha em $G_j$ um ponto $v_j$ mais próximo de $f(t_j)$, com desempate determinístico. Então $|v_j-f(t_j)|\le\delta/2$. As grades de paridades opostas são disjuntas, de modo que $v_{j+1}\ne v_j$. Como $N$ e $\delta$ são diádicos, o interpolante $g$ dos dados $(t_j,v_j)$ é uma montagem de átomos. Além disso,
\[
\norm{f-g}\le\norm{f-P_kf}+\norm{P_kf-g}
\le\frac M{8N^2}+\frac\delta2\le\varepsilon.
\]
Defina $\CC_\varepsilon$ como a imagem desse mapa de quantização. É exatamente o mesmo conjunto finito de interpolantes; resta apenas contá-lo usando a restrição de curvatura, em vez de contar todos os vetores da grade. Há $O(\varepsilon^{-1})$ escolhas para $v_0$. Dado $v_0$, o controle \eqref{eq:70} e os erros de quantização dão
\[
|v_1-v_0|\le C_0h+\delta,
\]
de modo que há $O(h/\delta+1)=O(\varepsilon^{-1/2})$ escolhas para $v_1$. Para $1\le j\le N-1$, a identidade de diferença segunda já usada acima e $Mh^2\le4\varepsilon$ implicam
\[
|v_{j+1}-2v_j+v_{j-1}|\le Mh^2+2\delta\le6\varepsilon.
\]
Fixados $v_{j-1}$ e $v_j$, o valor $v_{j+1}$ fica, portanto, confinado a um intervalo de comprimento $12\varepsilon$. Os pontos de $G_{j+1}$ têm passo $\delta>\varepsilon/2$, logo esse intervalo contém no máximo uma constante universal $K$ de candidatos; por exemplo, $K=25$ serve. Consequentemente,
\[
|\CC_\varepsilon|\le C\varepsilon^{-1}C'\varepsilon^{-1/2}K^{N-1},
\]
e, como $N=O(\varepsilon^{-1/2})$,
\[
\log_2|\CC_\varepsilon|\le O(\log(1/\varepsilon))+N\log_2K=O(\varepsilon^{-1/2}).
\]
Isso prova \eqref{eq:87} sem alterar os átomos, as grades alternadas ou a garantia de distorção.
\end{proof}

\subsection{Cobertura, bits, serialização, seleção e decodificação}
Seja $(X,d)$ um espaço métrico, $\FF\subset X$ totalmente limitado e $\varepsilon>0$.
\begin{definicao}[recursos de um codec]
Um $\varepsilon$-codec para $\FF$ é uma tripla $(W,\operatorname{Enc},\operatorname{Dec})$, na qual $W$ é um conjunto finito de palavras binárias, $\operatorname{Enc}:\FF\to W$ e $\operatorname{Dec}:W\to X$ satisfazem
\[
\sup_{f\in\FF}d\bigl(f,\operatorname{Dec}(\operatorname{Enc}(f))\bigr)\le\varepsilon.
\]
O decodificador é compartilhado e seu armazenamento não é contado no comprimento das palavras. Distinguem-se:
\begin{align*}
\mathcal H_\varepsilon(\FF)&=\log_2\mathcal N_\varepsilon(\FF)
&&\text{entropia métrica},\\
L_\varepsilon^{\rm dec}(\FF)&=
\inf_{\substack{\varepsilon\text{-codecs}\\\text{prefix-free}}}\max_{w\in W}|w|
&&\text{menor comprimento máximo decodificável},\\
L_\varepsilon^{\rm global}(\CC_\varepsilon)&=\lceil\log_2|\CC_\varepsilon|\rceil
&&\text{índice fixo em codebook compartilhado}.
\end{align*}
Para um codebook de montagens e uma convenção de serialização fixada, ponha
\[
L_{\rm folhas}(g\mid K)=L_{\rm arv}(g\mid K)+
\sum_{J\in\Lambda(g)}L_{\rm addr}(g|_J),
\qquad
L_\varepsilon^{\rm folhas}(\CC_\varepsilon\mid K)=
\max_{g\in\CC_\varepsilon}L_{\rm folhas}(g\mid K).
\]
Fixados um modelo de computação $\mathfrak M$, uma representação de entrada para $f$ e o codec, $T_\varepsilon^{\rm enc}$ é o custo de selecionar ou produzir a palavra e $T_\varepsilon^{\rm dec}$ é o custo de reconstruí-la.
\end{definicao}
O índice global nomeia uma entrada da tabela compartilhada; $L_\varepsilon^{\rm dec}$ mede comprimento de palavra. Os tempos são dados por $T_\varepsilon^{\rm enc}$ e $T_\varepsilon^{\rm dec}$.

\begin{proposicao}[limite informacional para representação decodificável]
Todo $\varepsilon$-codec de comprimento fixo $\ell$ para $\FF$ satisfaz
\eqn{89}{\ell\ge\lceil\mathcal H_\varepsilon(\FF)\rceil.}
A mesma desigualdade vale se $W$ é prefix-free e $\ell$ é o maior comprimento de palavra. Em particular,
\[
L_\varepsilon^{\rm dec}(\FF)\ge\lceil\mathcal H_\varepsilon(\FF)\rceil.
\]
Sob convenções que acrescentem ou removam apenas um cabeçalho fixo, a forma invariante é $\ell\ge\mathcal H_\varepsilon(\FF)-O(1)$.
\end{proposicao}
\begin{proof}
As reconstruções $\operatorname{Dec}(W)$ formam uma cobertura de $\FF$ por bolas de raio $\varepsilon$. Portanto,
\[
\mathcal N_\varepsilon(\FF)\le|\operatorname{Dec}(W)|\le|W|.
\]
No caso de comprimento fixo, $|W|\le2^\ell$. No caso prefix-free, a desigualdade de Kraft e $|w|\le\ell$ dão
\[
|W|2^{-\ell}\le\sum_{w\in W}2^{-|w|}\le1,
\]
e novamente $|W|\le2^\ell$. Tomar logaritmos prova \eqref{eq:89}. Essa separação entre entropia e complexidade de descrição é a que se usa em limites informacionais para esquemas quantizados \citep{bolcskei2019}.
\end{proof}
Aplicada ao teorema~7.18, a proposição fornece
\eqn{90}{\mathcal H_\varepsilon(\FF_{R,M})\le L_\varepsilon^{\rm dec}(\FF_{R,M})
\le L_\varepsilon^{\rm global}(\CC_\varepsilon)=O(\varepsilon^{-1/2}).}
A cota inferior $\Omega(\varepsilon^{-1/2})$ prova a otimalidade em ordem quanto à cardinalidade e aos bits do índice global. A serialização folha a folha da mesma construção segue o teorema~7.9 e pode usar $O(\log(1/\varepsilon))$ bits por folha, pois as alturas são polinomiais em $1/\varepsilon$.

% === FIM DA FONTE: secoes/07b_erro_entropia.tex ===


% === INICIO DA FONTE: secoes/07c_catalogos.tex ===
\subsection{Catálogos finitos, minimizadores e truncagem}
A existência de minimizadores depende da classe de busca. Na biblioteca infinita, o ínfimo pode deixar de ser atingido mesmo para uma única avaliação.
\begin{proposicao}
Para $x=1/4$ e $y=\sqrt2$,
\[
\inf\{|y-\ell_n(x)|:n\in\N_0,\ x\in I_n^x\}=0,
\]
mas nenhum átomo realiza o ínfimo.
\end{proposicao}
\begin{proof}
Para todo átomo, $a_n=y_{0,n}-s_nx_{0,n}\in\Q$ e $b_n=s_n\in\Q$; como $x=1/4$ é racional, $\ell_n(x)\in\Q$. Portanto, $\ell_n(x)\ne\sqrt2$.

Por outro lado, dado $\eta>0$, escolha $q\in\Q$ com $|q-\sqrt2|<\eta$. Use a célula $J_{2,1}=[1/4,1/2]$, de centro $x_0=3/8$, escolha $s=1$ e $y_0=q-s(x-x_0)=q+1/8\in\Q$. O teorema~3.6 fornece um átomo nessa célula com $\ell_n(1/4)=q$. Assim, os erros tendem a zero, mas nenhum é zero.
\end{proof}
Fixe uma raiz diádica $K=J_{k_0,m_0}$ e seja $\mathcal T_D(K)$ o conjunto das partições obtidas por árvores binárias diádicas completas de profundidade relativa no máximo $D$. Uma montagem representada é um par $\mathscr R=(\Lambda,n)$ formado por $\Lambda\in\mathcal T_D(K)$ e um endereço $n(J)$ para cada folha, com suporte exatamente $J$ e concordância dos valores nas fronteiras.

\begin{definicao}[catálogo truncado]
Para uma montagem representada, ponha
\[
\heightq_{\max}(\mathscr R)=
\max_{J\in\Lambda}\max\{2,k(J)+1,|m(J)|+1,\heightq(y_{0,J}),\heightq(s_J)\}.
\]
Se $H\ge H_K:=\max\{2,k_0+1,|m_0|+1\}$, defina
\eqn{91}{\RR^{\LL}_{D,H}(K)=
\{\mathscr R=(\Lambda,n):\Lambda\in\mathcal T_D(K),\ \heightq_{\max}(\mathscr R)\le H\},}
\eqn{92}{\Sigma^{\LL}_{D,H}(K)=
\{\Eval(\mathscr R):\mathscr R\in\RR^{\LL}_{D,H}(K)\}\subset\Asm(K).}
Representações distintas que avaliam na mesma função correspondem a um único elemento de $\Sigma^{\LL}_{D,H}(K)$.
\end{definicao}
Para fixar o desempate, serialize a árvore em pré-ordem com $0$ para nó interno e $1$ para folha e, depois, escreva os endereços $n^\sharp(J)+1$ das folhas em ordem espacial com o código gama de Elias. O esqueleto informa quantas palavras devem ser lidas, logo o fluxo $\ser^\sharp(\mathscr R)$ é autodelimitado. Ordene os fluxos por shortlex: primeiro pelo comprimento e depois lexicograficamente. Para cada função, retenha o menor fluxo entre suas representações no catálogo.

\begin{teorema}[finitude, existência e desempate canônico]
O catálogo $\Sigma^{\LL}_{D,H}(K)$ é finito e não vazio. Mais precisamente, se
\[
T_0=1,\qquad T_{d+1}=1+T_d^2,
\]
então
\eqn{93}{1\le|\Sigma^{\LL}_{D,H}(K)|\le|\RR^{\LL}_{D,H}(K)|
\le T_D|\LL_H|^{2^D}<\infty.}
Consequentemente, para qualquer funcional $\mathcal J:\Sigma^{\LL}_{D,H}(K)\to\R$, existe
\eqn{94}{g^\star\in\operatorname*{arg\,min}_{g\in\Sigma^{\LL}_{D,H}(K)}\mathcal J(g).}
Entre todos os minimizadores, a regra do menor fluxo shortlex seleciona uma única função e uma única representação canônica.
\end{teorema}
\begin{proof}
A recorrência de $T_D$ conta todas as formas ordenadas de árvore: ou a raiz é folha, ou seus dois filhos recebem independentemente uma forma de profundidade no máximo $D-1$. Cada árvore tem no máximo $2^D$ folhas e cada folha recebe um rótulo de $\LL_H$. Ignorar as restrições de suporte e continuidade só aumenta a contagem, o que prova a cota superior. A árvore reduzida à raiz pode ser rotulada, por exemplo, com $y_0=0$ e $s=1$, logo o catálogo não é vazio. A imagem de um conjunto finito por $\Eval$ é finita; qualquer função real em um conjunto finito não vazio atinge seu mínimo.

Resta o desempate. Há um número finito de representações de cada função, e a ordem shortlex é total; portanto existe um menor fluxo. A serialização é injetiva sobre representações, pois recupera a árvore e cada endereço de folha. Assim, dois candidatos distintos não compartilham o mesmo fluxo, e o menor minimizador é único.
\end{proof}

\begin{proposicao}[consistência sob truncagem]
Se $D\le D'$ e $H\le H'$, então
\eqn{95}{\Sigma^{\LL}_{D,H}(K)\subseteq\Sigma^{\LL}_{D',H'}(K).}
Para um funcional fixado na união dos catálogos, os valores mínimos são não crescentes em $D$ e $H$. Se o funcional possui um minimizador global nessa união, o minimizador global de menor fluxo shortlex e sua representação canônica pertencem a algum catálogo finito; a partir desse par de truncagens, o seletor canônico estabiliza.
\end{proposicao}
\begin{proof}
Uma representação admissível sob $(D,H)$ continua admissível quando os dois limites aumentam, o que prova \eqref{eq:95} e a monotonicidade dos mínimos. A ordem shortlex bem ordena as palavras binárias finitas. Logo, entre todas as representações de todos os minimizadores globais, existe uma menor; ela tem profundidade e altura finitas. Qualquer catálogo que já a contenha não pode ser superado posteriormente nem no valor do funcional nem no desempate.
\end{proof}

\begin{algoritmo}[enumeração exata do catálogo]
\leavevmode\par\nobreak\noindent
Para parâmetros pequenos, gere recursivamente todas as formas de $\mathcal T_D(K)$; em cada folha, liste somente os átomos de $\LL_H$ com o suporte exigido; ao unir duas subárvores, descarte pares cujos valores na fronteira diferem. Reduza representações funcionalmente idênticas ao menor fluxo $\ser^\sharp$, avalie $\mathcal J$ em aritmética racional exata e aplique o desempate shortlex.
\end{algoritmo}
O algoritmo realiza o mínimo do teorema~7.23 por enumeração do catálogo. A figura~\ref{fig:circuito} reúne as etapas da construção.
\\begin{figure}[!htbp]\\centering
 \\emph{Diagrama de circuito omitido nesta versao para formalizacao.}
 \\caption{Circuito das etapas da construcao.}\\label{fig:circuito}
 \\end{figure}
\clearpage

% === FIM DA FONTE: secoes/07c_catalogos.tex ===


% === FIM DA FONTE: secoes/07_codificacao.tex ===


% === INICIO DA FONTE: secoes/08_conclusao.tex ===
\section{Conclusão}
A biblioteca $\LL$ reúne segmentos afins racionais com endereço único e suporte diádico recuperável. A bijeção com jatos racionais produz montagens contínuas densas em $C(K)$ e uma representação ReLU rasa exata. Na aproximação uniforme de funções $C^2$, a racionalização preserva a constante $1/8$ do erro interpolatório.

O refinamento adaptativo combina majorantes intervalares da segunda derivada e racio\-na\-li\-za\-ção nodal para obter erro até $\varepsilon_I+\varepsilon_R$, com término e número de folhas controlados. Em $\FF_{R,M}$, a construção conjunta fornece erro $O(N^{-2})$, altura $O(N^2)$ e descrição $O(N\log N)$.

Sob o calibre fixado, a contagem $\Theta(H^6)$ determina o limite minimax do endereçamento, realizado pela ordem por altura. Os formatos estruturais usam $O(N\log H)$ bits, e o codebook de montagens realiza a ordem de entropia $\Theta(\varepsilon^{-1/2})$. A truncagem simultânea de profundidade e altura fecha a construção com catálogos finitos, existência de minimizadores e uma representação canônica definida por shortlex.

% === FIM DA FONTE: secoes/08_conclusao.tex ===


% === INICIO DA FONTE: secoes/A_notacao.tex ===
\section{Quadro de notação}
\footnotesize
\renewcommand{\arraystretch}{1.1}
\begin{longtable}{@{}p{.28\linewidth}p{.68\linewidth}@{}}
\toprule Símbolo&Significado\\\midrule
\endfirsthead
\toprule Símbolo&Significado\\\midrule
\endhead
\bottomrule\endfoot
$\pi,\pi^{-1}$&Pareamento de Cantor e sua inversa.\\
$Z,Q,Q^*$&Enumerações de $\Z$, $\Q$ e $\Q^*$.\\
$\heightq(p/q)$&Altura $\max\{|p|,q\}$ da fração reduzida.\\
$\Q_H$&Racionais de altura no máximo $H$.\\
$\D,\D_{\rm mid}$&Diádicos e diádicos não inteiros.\\
$J_{k,m},J_{k,m}^{\rm own}$&Célula fechada e célula semiaberta de posse.\\
$x_0,y_0,s$&Centro, valor ancorado e inclinação de um átomo.\\
$\Gamma,\LL,\LL_H$&Gerador aritmético, biblioteca e truncagem por altura.\\
$I_n^x$&Suporte espacial recuperado da interface global.\\
$\Asm(K)$&Montagens contínuas finitas no domínio $K$.\\
$\omega(f;\delta)$&Módulo de continuidade de $f$.\\
$\rho(z)=z_+$&Ativação ReLU.\\
$P_kF,P_Jf,P_\Lambda F$&Interpolantes na malha uniforme, em uma célula e em uma partição.\\
$\EE,\mathbb I_{\Q}$&Gramática de expressões e intervalos de extremos racionais.\\
$P,d$&Precisão transcendental e profundidade de subdivisão intervalar.\\
$M_{\rm int}(J),M(J)$&Majorantes validados de $\|F''\|_{\infty,J}$.\\
$\eta_{\rm cert}(J)$&Certificado local $|J|^2M(J)/8$.\\
$\varepsilon_I,\varepsilon_R$&Tolerâncias de interpolação e racionalização.\\
$H$&Calibre atômico definido em \eqref{eq:46}.\\
$\Gamma^\sharp,n^\sharp$&Enumeração por altura e endereço nessa ordem.\\
$\gamma,C_{\Z},C_{\rm rat},C_{\rm dyad}$&Código gama e códigos de inteiros, racionais e diádicos.\\
$L_{\rm atom},L_{\rm ReLU,str},L_{\rm ReLU,rat}$&Comprimentos dos três formatos de serialização.\\
$\Asm_{N,H}(K)$&Montagens com até $N$ folhas e altura atômica até $H$.\\
$E^{\LL}_{N,H}(f)_\infty$&Melhor erro uniforme sob os dois orçamentos.\\
$H^{\LL}_{\varepsilon,N}(f)$&Menor altura suficiente para erro até $\varepsilon$ com até $N$ folhas.\\
$\FF_{R,M}$&Classe $\{f\in C^2[0,1]:\|f\|_\infty\le R,\ \|f''\|_\infty\le M\}$.\\
$\mathcal N_\varepsilon,\mathcal H_\varepsilon$&Número de cobertura e entropia métrica.\\
$\CC_\varepsilon$&Codebook finito de montagens com distorção até $\varepsilon$.\\
$L_\varepsilon^{\rm dec},L_\varepsilon^{\rm global},L_\varepsilon^{\rm folhas}$&Comprimentos decodificável, do índice global e da serialização por folhas.\\
$T_\varepsilon^{\rm enc},T_\varepsilon^{\rm dec}$&Custos de seleção da palavra e reconstrução no modelo fixado.\\
$K=J_{k_0,m_0},\mathcal T_D(K)$&Raiz diádica e partições por árvores de profundidade relativa até $D$.\\
$\RR^{\LL}_{D,H}(K),\Sigma^{\LL}_{D,H}(K)$&Catálogos de representações e de funções sob profundidade e altura.\\
$\ser^\sharp$&Serialização autodelimitada usada no desempate shortlex.\\
$\mathcal J,g^\star$&Funcional de seleção e minimizador no catálogo.\\
\end{longtable}
\normalsize

% === FIM DA FONTE: secoes/A_notacao.tex ===
\end{document}
